\begin{savequote}
\includegraphics[width=3.5cm]{./images/pallinger-christoph-aass/Anton_32986_v2-AASS-0112mm.jpg}

Bild: Christoph Pallinger
\end{savequote}
\chapter
[\CO{[SAN]} Angewandte Sanitätshilfe]
{\CC{[SAN]} Angewandte Sanitätshilfe}
% \AasRsN{RS~XIII}
\label{chp:SAN}

\ChapterIntro
{%
Dieses Kapitel behandelt diverse allgemeine und
spezielle Maßnahmen der Sanitätshilfe. 
}
{%
\begin{description}
  \item[Maintainer:] Diverse
  \item[Autoren:] Diverse
  \item[Reviewer:] --
  \item[Version:] {\DocVersion} ({\DocVersionInfo})
  \item[SHA1:] {\DocGitCommit}
\end{description}

{\TxTeilkapitelAASS}
}

% \author{Gabriel, Hochreiter, Motal}

% \minitoc
% \includegraphics[angle=270,width=\linewidth]{./images/leitsystem/Leitschema-PAM_SAN.pdf}

\clearpage % TODO temp clearpage
\Layout{0}
{\TxQuerverweise}
{~\par\noindent\includegraphics[angle=270,width=\linewidth]{./images/leitsystem/Leitschema-PAM_SAN.pdf}}
{\begin{itemize}
  \item \EE{Beschreibung der Medizinprodukte}: \CO{[MP1, MP2]} 
  (\ref{chp:MP1}, \ref{chp:MP2}).
  \item \EE{Untersuchungen}: \CO{[BAU]} (\ref{chp:BAU}).
  \item \EE{Sanitätshilfemaßnahmen}:  \CO{[SAN]} (\ref{chp:SAN}).
%   \item[Anamnese und Patientengespräch:]  Kapitel [AUP] (\ref{chp:AUP}).
\end{itemize}}
{}
{}
{}
% 
% \Layout{2}{Inhalt}{%
% Dieses Kapitel behandelt die Durchführung von Maßnahmen der Sanitätshilfe.
% }{}{}{}{}
% 
% 
% \Layout{2}{{\TxQuerverweise}}{%
% \begin{compactdesc}
%   \item[Beschreibung der Medizinprodukte:] Kapitel [MP1, MP2] 
%   (\ref{chp:MP1}, \ref{chp:MP2}).
%   \item[Untersuchungen:] Kapitel [BAU] (\ref{chp:BAU}).
%   \item[Sanitätshilfemaßnahmen:]  Kapitel [SAN] (\ref{chp:SAN}).
% %   \item[Anamnese und Patientengespräch:]  Kapitel [AUP] (\ref{chp:AUP}).
% \end{compactdesc}
% }{}{}{}{}


% \clearpage
\section{Allgemeine Tätigkeiten}

\section{Absaugung}

\label{topic:absaugkatheter}
\label{topic:absaugung-taetigkeit}
\index{Absaugung}
\index{Absauggeräte}
\index{Absauggeräte!Accuvac{\texttrademark}}
\index{Absauggeräte!Laerdal Suction Unit{\texttrademark}}
\index{Accuvac{\texttrademark}}
\index{Laerdal Suction Unit{\texttrademark}}
\index{Absaugkatheter}
\label{topic:absaugung}
\label{topic:absauggeraete}

\Layout{2}{{\TxQuerverweise}}{%
\begin{inparaitem}
  \item Absaugkatheter: \prettyref{topic:absaugkatheter};
  \item Absaugtechnik: \prettyref{topic:absaugung-taetigkeit}
\end{inparaitem}
}{}{}{}{}

\Layout{2}{Einleitung}
{Blut oder Schleim kann vom Fachpersonal mittels
Absauggeräten aus dem Mund-Rachen-Raum abgesaugt werden. Es
gibt unterschiedliche Modelle; grundsätzlich wird nach der
Betriebsart unterschieden zwischen \EE{elektrischer
Absaugeinheit}, \EE{manueller Absaugpumpe} und als
Sonderfall der \EE{Oro-Sauger} oder Schleimabsauger.

Ein \EE{Absaugkatheter} ist jenes Zubehör zum Absauggerät,
mit welchem die Sekrete aus Mund oder Nase abgesaugt werden
können. Es ist ein flexibler, ca. 50\CM* langer Schlauch,
der direkt an die Absaugpumpe angeschlossen wird (siehe
Abbildung in \prettyref{tab:absaugvorrichtungen}). Die
meisten Absaugkatheter haben auf Patientenseite eine kleine
Öffnung, die das Festsaugen an den Patienten verhindert. Die
Absaugkatheter sind steril verpackt und müssen auch
dementsprechend behandelt werden. Die farbige Kennung gibt
den Durchmesser des Absaugkatheters an. \Z{Die Farbcodierung
kann \prettyref{tab:absaugkatheter} entnommen werden.}}
{}
{}
{}
{}



\Layout{57}{}
{}{}
{\Z{%
\begin{tabulary}{\linewidth}{LCCCCCCCCC}
\noindent\TabHeading{Farbe}
&
\noindent\TabHeading{grau}
&
\noindent\TabHeading{hellgrün}
&
\noindent\TabHeading{blau}
&
\noindent\TabHeading{schwarz}
&
\noindent\TabHeading{weiß}
&
\noindent\TabHeading{grün}
&
\noindent\TabHeading{orange}
&
\noindent\TabHeading{rot}
&
\noindent\TabHeading{gelb}
\\\addlinespace\midrule
Durchmesser in mm
& 1,7 & 2,0 & 2,7 & 3,3 & 4,1 & 4,7 & 5,3 & 6,1 & 6,6
\\\addlinespace
Größe in Charrière 
& CH\,5 & CH\,6 & CH\,8 & CH\,10 & CH\,12 & CH\,14 & CH\,16
& CH\,18 & CH\,20 
\\\addlinespace\bottomrule
\end{tabulary}
}
}
{Farbcodierung: Absaugkatheter\label{tab:absaugkatheter}}{}

\opt{ORGASBFLO}{\ASBFLO{Beim ASB Floridsdorf-Donaustadt werden folgende Geräte
eingesetzt:

\Layout{57}{}{}{}
{%
\begin{tabulary}{\linewidth}{LLLL}
\TabHeading{Hersteller}
&
\TabHeading{Produktfamilie}
&
\TabHeading{Typ}
&
\TabHeading{Bemerkungen}
\\\addlinespace\midrule
Weinmann
&
\noindent\T{Accuvac{\texttrademark}}
&
Rescue
&
rote Deckplatte
\\\addlinespace
&
&
Basic
&
grüne Deckplatte
\\\addlinespace
Laerdal
& 
Laerdal Suction Unit{\texttrademark}
&
&
auslaufend
\\\addlinespace
div. Hersteller
&
Handabsaugpumpe
&
RES-Q-VAC
&
\\\addlinespace
div. Hersteller
&
Orosauger
&
&
\\\addlinespace\bottomrule
\end{tabulary}
}
{ASB921: Eingesetzte Absauggeräte.}
{}
}}%Ende opt

\Layout{57}{}{}{}
{
\begin{tabulary}{\linewidth}{LLL}
\TabHeading{Elektrische Absaugeinheit}

Älteres Modell von der Fa. Laerdal (Laerdal Suction Unit{\texttrademark})

Neueres Modell von der Fa. Weinmann (Accuvac{\texttrademark})
&
\vspace{0.1pt}
\includegraphics[width=\linewidth]{./images/pallinger-christoph-aass/LSU_32831_export_00800.jpg}
&
\vspace{0.1pt}
\includegraphics[width=\linewidth]{./images/pallinger-christoph-aass/Accuvac_32820-AASS-0112mm.jpg}
\\\midrule
\TabHeading{Manuelle Absaugpumpen}

Handabsaugpumpe

Fußabsaugpumpe
&
\vspace{0.1pt}
\includegraphics[width=\linewidth]{./images/geraete/handabsaugung.jpg}
&
\vspace{0.1pt}
\includegraphics[width=\linewidth]{./images/geraete/fussabsaugung1.jpg}
\\\midrule
\TabHeading{Oro-Sauger}

Schleimabsauger für Neugeborene
&
\vspace{0.1pt}
\includegraphics[width=\linewidth]{./images/geraete/orosauger.jpg}
\\\midrule
\TabHeading{Absaugkatheter}

zum Einführen in Mund oder Nase
&
\vspace{0.1pt}
 \includegraphics[width=4cm]{./images/geraete/absaugkatheter.jpg} \\\bottomrule
\end{tabulary}
}
{Verschiedene Absaugvorrichtungen.\label{tab:absaugvorrichtungen}}
{}

% \begin{table}[H]
% \begin{WidePar}
% \Captionof{table}{}
% 
% 
% \end{WidePar}
% \end{table}

\Layout{0}{Funktionskontrolle}
{%
Jede Funktionskontrolle umfasst die Überprüfung der
\EE{Vollständigkeit} des Materials, \EE{Saugfunktion} bei
maximaler und minimaler Leistung und \EE{Dichtheit}. Sie muss
zumindest bei Dienstbeginn durchgeführt werden.

Bei Geräten, welche über Akkus verfügen, muss die
\EE{Akku-Ladung} und Saugfunktionsprüfung \E{ohne Anschluss
an eine externe Stromquelle} erfolgen, d.h. das Gerät muss
für die Funktionsprüfung von einer Ladehalterung abgenommen
werden. Die Ladestandsanzeige und evtl. vorhandene
\EE{Warnleuchten} müssen kontrolliert werden.

Die vom jeweiligen Hersteller vorgegebenen,
typenspezifischen Anweisungen müssen beachtet werden. 
}{
\begin{itemize}
   \item Prüfung von
   \begin{itemize}
      \item Vollständigkeit
      \item Saugleistung
      \item Dichtheit
      \item Ggfs. Akku-Ladung
      \begin{itemize}
         \item Ohne externe Stromquelle!
      \end{itemize}
      \item Warnleuchten
   \end{itemize}
   \item Herstellerhinweise beachten!
\end{itemize}
}{}{}{}



\Layout{37}{Weinmann AccuVac{\protect\texttrademark}}
{}
{}
{./images/pallinger-christoph-aass/Accuvac_32820-AASS-0176mm.jpg}
{Weinmann AccuVac{\protect\texttrademark}}
{Ch. Pallinger}




\Layout{0}{Absaugbereitschaft}
{%
\index{Absaugbereitschaft}
Wenn \EE{Absaugbereitschaft} hergestellt werden muss, wird
aus hygienischen Gründen nur der farbige Trichter des
Katheters ausgepackt und an den Konnektor bzw. den
\T{Fingertip} der Absaugpumpe angeschlossen. Der restliche
Teil des Katheters bleibt so lange in der Verpackung bis mit
der Absaugung begonnen
wird.\bigskip}{
\begin{itemize}
  \item Verpackung nur am farbigen Ende des Katheters öffnen
  \item An Absaugpumpe anschließen
  \item Katheter bleibt in Hülle
\end{itemize}
}
{}
{}
{}

\Layout{0}{Absaugtechnik}{%
Damit der
Absaugkatheter nicht zu tief eingeführt wird, kann die
einzuführende Katheterlänge durch Abmessen der Distanz
zwischen \EE{Mundwinkel und Ohrläppchen} bzw. zwischen \EE{Nasenspitze
und Ohrläppchen} bestimmt werden. Es darf \EE{maximal auf Sicht} abgesaugt
werden!

Der Katheter wird \EE{ohne Sog eingeführt}, dazu wird entweder der Katheter
abgeknickt  oder der Fingertip offen gelassen. Unter kreisenden Bewegungen wird
dann auf Sicht abgesaugt und der Katheter herausgezogen. Wenn sich der Katheter
festsaugt, muss der Fingertip geöffnet oder der Katheter abgesteckt werden.
Der Vorgang wird nach Bedarf wiederholt.

\Cave{Es darf nur auf Sicht abgesaugt werden!}
}{
\begin{itemize}
  \item Tiefe:
  \begin{itemize}
    \item Max. auf Sicht.
    \item Mund: Mundwinkel -- Ohrläppchen.
    \item Nase: Nasenspitze -- Ohrläppchen.
  \end{itemize}
  \item Ohne Sog einführen
  \begin{itemize}
    \item Katheter abknicken, Fingertip offen
  \end{itemize}
  \item Unter kreisenden Bewegungen saugen und herausziehen
\end{itemize}
}{}{}{}

\begin{table}[H]
\FormatTable
\Captionof{table}{Bilderserie: Absaugtechnik}
\begin{tabularx}{\linewidth}{XX}
\includegraphics[width=\linewidth]{./images/pallinger-christoph-aass/Anton_32999_export.jpg}
\Captionof{figure}{Katheterlänge bei Absaugung des Mundes \SourcePicture{Ch. Pallinger}}
&
\includegraphics[width=\linewidth]{./images/pallinger-christoph-aass/Anton_32990_export.jpg}
\Captionof{figure}{Katheterlänge bei Absaugung durch die Nase \SourcePicture{Ch. Pallinger}}
\end{tabularx}
\end{table}





\Layout{0}{\TxGefahren}{%
In der Rachenhinterwand verläuft ein Nerv, welcher Einfluss auf den Kreislauf
hat (Vagus-Nerv). Wenn dieser gereizt wird kann es zu gefährlichen
\EE{Kreislaufstörungen} (Bradykardie, Hypotension) kommen. Der Vagus-Nerv ist
Teil des \E{Parasympathikus} und des vegetativen Nervensystems
(\prettyref{topic:parasymapthikus}). 
}{
\begin{itemize}
  \item Reizung des Vagus-Nerv
  \item Gefährliche Kreislaufstörungen
\end{itemize}
}{}{}{}


\section{Sauerstoff -- \ce{O2}}
\label{topic:sauerstoff}
\label{topic:o2}
\index{Sauerstoff}
\index{O2@\ce{O2}}

\Definition{\T{Sauerstoff} ist ein
lebenswichtiges Gas welches zu 21\PC* in unserer Atemluft
vorkommt und außerdem ein essentielles
Notfall"`medikament"'.
Das chemische Symbol lautet \T{\ce{O2}}.}

Sauerstoff darf und muss bei Bedarf vom Sanitäter
verabreicht werden. Die frühzeitige Verabreichung ist wichtig für die Prognose des
Patienten!

% \begin{center}
% \includegraphics[width=4.522cm,height=5.628cm]{./images/rippenspreitzer-02-fett.png}
% \Captionof{figure}{NEIN! \SourcePicture{Rippenspreitzer.de}}
% 
% \end{center}

\subsection{Technische Grundlagen}
% TODO : Layout-Anpassung
\TODO{GABS}{2010-mm-dd}{Layout-Anpassung notwendig.}
{Dieser Teil entspricht nicht dem Standard-Layout!}
\TODO{GABS}{2010-08-04}{Bild: Sauerstofffl., Druckminderer fehlt.}{}
\TODO{KOCH}{2010-11-06}{Kennzeichnung O2-Flasche ohne "`N"'?}{}
%
Sauerstoff ist ein farb- und geruchsloses Gas und wird in
Druckflaschen gelagert.

\Layout{0}{Kennzeichnung}
{Die \EE{Farbe} von Gasflaschen ist
normiert\ZFootnote{Normen zur Kennzeichnung von Gasflaschen: EN~1089-3, in der jeweiligen nationalen
Umsetzung (z.\,B. DIN EN~1089-3, ÖNORM EN~1089-3
(Verbindlichkeit gem. Versandbehälterverordnung~2002 (VBV~2002)), \ldots).}.
Flaschen mit Sauerstoffbefüllung werden an der
Flaschenschulter mit einem \EE{weißem Farbring} gekennzeichnet.
Gase für den medizinischen Gebrauch werden zusätzlich
durch einen weißen Flaschenkörper gekennzeichnet.
Gasflaschen mit \EE{medizinischem Sauerstoff} sind daher
\EE{vollständig weiß}\ZFootnote{Gasflaschen für
\E{technischen} Sauerstoff haben zwar eine weiße
Flaschenschulter, jedoch einen \E{blauen Flaschenkörper}.
Technischer Sauerstoff darf Patienten nicht verabreicht
werden.}.
\cite{AUT-LAW-VBV2002-1,Igv2010} 

Weiters ist auf Druckgasflaschen ein
\EE{Gefahrengutaufkleber} angebracht.

\Fazit{Gasflaschen mit \EE{\E{medizinischem} Sauerstoff}
sind
\EE{vollständig weiß}.}

\Z{Die Farbcodierung wurde in der Zeit von
\VonBis{1998}{2006}\ZFootnote{Die Umstellung auf das neue
Farbsystem gemäß ÖNORM EN 1089-3 musste
bei allen anderen Gasen bis zum \E{30. Juni 2006}
abgeschlossen sein (\S\,26 Abs.\,3 VBV 2002
\cite{AUT-LAW-VBV2002-1}).} umgestellt. In der Übergangszeit wurden Flaschen mit der neuen Codierung mittels Aufbringung des
Buchstabens "`\EE{\sffamily N}"' ("`neue Kennzeichnung"')
gekennzeichnet. Dies hat \E{nichts} mit dem
chemischen Symbol für Stickstoff (\ce{N}) zu tun!} 

\Commentary[Gasflaschen / Farbcodierung / Umstellung]
{\S\,26 Abs.\,2 VBV 2002:
Bei Sauerstoff, Distickstoffmonoxid (Lachgas) und
oxidierenden Gasen erfolgt die Umstellung vom Farbsystem
gemäß \S\,15 Abs. 1 Z 2 lit. a für die industrielle Anwendung
ab dem 1. Jänner 2002 und bei solchen für medizinische
Zwecke ab dem 1. Jänner 2004.
\ParBugfix
Abs. 3: Die Umstellung auf das neue Farbsystem gemäß ÖNORM
EN 1089-3 hatte bei Acetylen bis zum 31. Dezember 2001 stattzufinden,
bei allen anderen Gasen muss sie bis zum 30. Juni 2006
abgeschlossen sein.

\cite{AUT-LAW-VBV2002-1}
}
\Commentary[Gasflaschen / Farbcodierung / EN~1089-3]
{Nach: \cite{Igv2010}:

Die verbindliche Kennzeichnung des Gaseinhalts erfolgt auf dem
Gefahrgutaufkleber.
Die Farbkennzeichnung ist nur für die Flaschenschulter festgelegt, 
 außer bei medizinischen Gasen. In diesem Fall ist der zylindrische Teil weiß.
Die Farbe des zylindrischen Flaschenmantels ist in der Norm (bis auf
medizinischen Gase) nicht festgelegt. 
}

}
{
\begin{itemize}
  \item Medizinischer Sauerstoff: Vollständig weiß
  \item Gefahrengutaufkleber
\end{itemize}
}
{}
{}
{}


\Layout{0}{Eigenschaften der Druckflaschen}
{Die Druckflaschen haben unterschiedliche Füllgrößen
(Volumina), gängige Flaschengrößen sind 0,5, 1, 1,5, 2, 5
und
10\Liter*. Der Sauerstoff wird unter Druck gelagert, wodurch
die Menge des zur Verfügung stehenden Sauerstoffs um ein Vielfaches
gesteigert wird. Der Druck einer gefüllten Flasche beträgt
normalerweise ca. 200\Bar*. 

Jede Druckflasche verfügt über ein \T{Hauptventil} und
ein Standardgewinde. An diesem Gewinde ist entweder ein
weiterführender Druckschlauch oder der \E{Druckminderer}
einer \E{Berieselungseinheit} angeschlossen. }
{
\begin{itemize}
  \item Unterschiedliche Füllgrößen
  \item Druck ca. 200\Bar*
  \item Hauptventil mit Standardgewinde
  \item Daran Druckminderer/Berieselungseinheit
  angeschlossen
\end{itemize}
}
{}
{}
{}


\Layout{0}{Berieselungseinheit}
{\index{Berieselungseinheit}
Die Berieselungseinheit besteht im Wesentlichen aus einem \T{Druckminderer}. Der
Druckminderer hat ein \I{Abgabeventil}, welches den \E{Druck drosselt} und eine
Abgabe an den Patienten ermöglicht. Je nach Bauart kann der Druckminderer über 
ein \E{Regelventil}, ein \E{Flowmeter}\index{Flowmeter} und ein \E{Manometer}
verfügen.  Mit dem \T{Regelventil} kann man die Durchfluss- bzw. Abgaberate
ein\-stellen,  das \T{Flowmeter} (Durchflussanzeige) zeigt die aktuelle
Durchflussrate in  Liter pro Minute (\LPerMin) an. Das \T{Manometer}  zeigt den
Druck in der Sauerstoffflasche an, wenn das Hauptventil geöffnet ist.}
{
\begin{itemize}
  \item Druckminderer
  \item Regelventil
  \item Flowmeter \Sign{→} \LPerMin
  \item Manometer \Sign{→} \Bar
\end{itemize}
}
{}
{}
{}






\subsection{Umgang mit Sauerstoff und Druckflaschen}

\Layout{0}{Besondere \Ca{Gefahrenquellen}}
{%
Sauerstoff ist \EE{brandfördernd} und zusammen mit Fett
sogar \Ca{selbstentzündlich}! Neben der Brand- und
Explosionsgefahr birgt der große Druck unter dem die Flaschen
stehen weitere Gefahren: Sollte ein Ventil abbrechen, kann die
Flasche eine enorme Kraft entwickeln und zu einer Rakete
werden. Die Kraft einer mit 200\Bar* gefüllten Flasche reicht
aus, um Mauern zu durchbrechen!


}{
\Ca{← Siehe Fließtext!}
}{}{}{}

\Layout{0}{Beenden der \ce{O2}-Gabe und Wechsel von Druckflaschen}
{Nach der Verwendung des Sauerstoffsystems wird das
Hauptventil geschlossen und der \EE{Überdruck} aus den
Druckleitungen bzw. dem Druckminderer \EE{abgelassen}. Erst
danach wird das Abgabeventil ebenfalls verschlossen.

Gleiches gilt für den Wechsel einer Sauerstoffflasche: Es
muss -- zur eigenen Sicherheit -- sichergestellt sein, dass
\EE{im Entnahmeschenkel \Ca{kein Überdruck}} herrscht. Vor
dem Abmontieren des Druckminderers oder einer Druckleitung
muss dies auf alle Fälle \EE{extra am Manometer kontrolliert
werden}! Das \EE{Abgabeventil} muss beim Abmontieren
\EE{geöffnet} sein.



}{
\Ca{← Siehe Fließtext!}
}{}{}{}



\Layout{2}{\Ca{Sicherheitshinweise}}
{\Cave{\Ca{Achtung! Vorsicht beim Hantieren -- Zerknall- und
Explosionsgefahr!}
\begin{enumerate}
    \item Sauerstoff ist zusammen mit Fett selbstentzündlich!
    \begin{enumerate}
        \item Die Armaturen und Gewinde sind fettfrei zu
        halten!
        \item Ventile dürfen nicht mit fetten oder öligen Händen
        bedient werden! 
    \end{enumerate}
    \item Das Hantieren mit Feuer und offenem
    Licht ist in der Nähe von Sauerstoffflaschen verboten!
    \item Die Druckflaschen sind immer zu sichern und vor
    Umfallen, etc. zu schützen!
\end{enumerate}

\Ca{Lebensgefahr!} Vor der Abmontage eines
Druckminderers oder einer Druckleitung muss die
Überdruckfreiheit geprüft und das Abgabeventil geöffnet
werden!

\Fazit{Die Sauerstoffflaschen sind pfleglich zu behandeln!}
}}
{}
{}
{}
{}

\subsection{Heiteres Rechnen mit Sauerstoff}

Mit Sauerstoff lässt sich gut rechnen ...

Die insgesamt verfügbare Menge wird errechnet, indem man
die Füllgröße der Flasche (in Liter) mit dem
Flaschendruck (in Bar) multipliziert. Wenn man wissen möchte, wie lange die vorhandene Sauerstoffmenge bei
einer bestimmten Abgaberate reicht, dividiert man einfach
die gesamt verfügbare Menge durch die Abgaberate. 

\Formula{
\begin{displaymath}
\mbox{Verfügbares Volumen in \Liter} = {\begin{array}{c}
\mbox{Volumen der Flasche in \Liter}\\
\mbox{(Füllgröße)} \end{array}}
\times 
\begin{array}{c}{\mbox{Druck in bar}}\\
\mbox{(Flaschendruck)}
\end{array}
\end{displaymath}

\begin{displaymath}
\mbox{Zeit} = \frac{\mbox{Verfügbare
Menge}}{\mbox{Abgaberate}}
\end{displaymath}
}


%oder kurz: 

%\begin{equation}
%l = {Volumen} \times {Druck}
%\end{equation}



\Ca{Achtung!} In der Flasche sollten min. \EE{10\Bar*
Restdruck} verbleiben und es ist eine \EE{ausreichende
Zeitreserve} einzuplanen!  

\Fazit{Es ist ein \EE{Restdruck} abzuziehen und eine
ausreichende \EE{Zeitreserve} einzuplanen!}



%oder kurz: 

%\begin{equation}
%min = \frac{l}{l / min}
%\end{equation}


\minisec{Beispiel}

\begin{quote}
Man hat eine Flasche mit einer Füllgröße von 10\Liter*. Diese
steht unter einem Druck von 150\Bar*: Daraus ergibt
sich, dass in der Flasche 1\,500\Liter* \ce{O2} sind.

Der Patient benötigt 5\,\unitfrac{\Liter}{\Min} \ce{O2} : 

Der Inhalt würde theoretisch für 300 Minuten reichen. Der
Restdruck sowie eine Zeitreserve sind dabei
aber noch \E{nicht} berücksichtigt!
\end{quote}

\minisec{Beispiel}

\begin{quote}
Gegeben ist: 
\begin{itemize}
  \item \ce{O2}-Flasche mit Füllgröße von 10\Liter*, Druck von
160\Bar*
\item Patient: \ce{O2}-Bedarf von 5\LPerMin*
\item Fahrtdauer laut Routenplaner:
4\Hour*~58\Min*
\end{itemize}

Nach Abzug des Restdruckes ergibt sich ein nutzbarer Druck
von 150\Bar* und folglich eine nutzbare \ce{O2}-Menge von
1\,500\Liter*. Der Patient könnte daher für 300\Min* mit
Sauerstoff versorgt werden. Die (geschätzte) Fahrtdauer
von 4\Hour*~58\Min* entspricht 298\Min*, es ergibt sich
eine Zeitreserve von 2\Min*.

Da die Zeitreserve von 2\Min* zu kurz ist, kann der
Transport mit dieser Ausstattung \E{nicht} durchgeführt
werden. 
\end{quote}

\subsection{Verabreichung von Sauerstoff}
\index{Sauerstoff!Verabreichung}


\Layout{60}{Einleitung}
{%
Die Verabreichung von Sauerstoff ist eine wichtige und häufige
Maßnahme. Sauerstoff kann sowohl einem selbst atmenden Patienten
zusätzlich gegeben werde (\T{Berieselung}), aber auch bei der
\EE{Beatmung} zugeführt werden. 
}
{}
{./images/motal-aass/o2_4765-00800.jpg}
{}
{Motal}


\Fazit{Unterscheide Berieselung (erhaltene Eigenatmung) und Beatmung!}

\subsubsection{Berieselung}
\index{Berieselung}

Die \ce{O2}-Berieselung kann \EE{nur bei vorhandener
Spontanatmung} durchgeführt werden!

% \marginpar{\T{Sauerstoff\-brille}}
% \marginpar{\T{Sauerstoff\-maske}}
% \marginpar{\T{Sauerstoff\-maske + Reservoir}}

% zwei kurze, in die Nasenöffnungen  eingeführte Kunst\-stoff\-stutzen
% 
% bei ca. 4\,l O2 /min bis 40\% \ce{O2}  in Atemgas  möglich
% 
% 
% Bei behinderter Nasenatmung kann die  Sauerstoffbrille im Mundwinkel
% plaziert werden. 



\begin{table}[H]
\FontTableBase
\begin{WidePar}
\Captionof{table}{Verabreichungsarten von Sauerstoff.}

\begin{tabulary}{\linewidth}[H]{CLLL}
\TabHeading{Was}
 &
\TabHeading{Wann}
 &
\TabHeading{Vorteile}
&
\TabHeading{Nachteile}
\\\addlinespace\toprule\addlinespace 
\TabSub{Sauerstoff\-brille}
 &
Spontanatmung

\ce{O2}-Flow {\textless}~5\LPerMin*
 &
gute Toleranz

fehlendes Engegefühl

Sprechen \& Husten möglich
 &
ungenaue Dosierung

Austrocknung  der Schleimhäute 
\\\addlinespace\midrule\addlinespace
\TabSub{Sauerstoff\-maske}
 &
Spontanatmung

\ce{O2}-Flow ab 5\LPerMin*
 &
höhere \ce{O2}-Konzentration
 &
\ce{CO2}-Rückatmung  wenn \ce{O2}-Durchfluss
{\textless}~5\LPerMin*!) 

für Patienten evtl. unangenehm bzw. gewöhnungsbedürftig
\\\addlinespace\midrule\addlinespace
\TabSub{Sauerstoff\-maske mit Reservoir}
 &
Spontanatmung

\ce{O2}-Flow ab 5\LPerMin*
 &
Noch höhere \ce{O2}-Konzentration
 &
Wie oben.

Reservoir muss zuerst händisch befüllt werden.
\\\addlinespace\bottomrule
\end{tabulary}
\end{WidePar}
\end{table}

\begin{table}[H]

\begin{WidePar}
\Captionof{table}{Bilderserie: Mittel zur Verabreichung von Sauerstoff.}
\begin{tabularx}{\linewidth}{XXX}
\includegraphics[width=\linewidth]{./images/pallinger-christoph-aass/Anton_32988_v2-AASS-0112mm.jpg}
\Captionof{figure}{Sauerstoffbrille. \SourcePicture{Ch. Pallinger}}
&
\includegraphics[width=\linewidth]{./images/pallinger-christoph-aass/Anton_32987_v2-AASS-0112mm.jpg}
\Captionof{figure}{Sauerstoffmaske. \SourcePicture{Ch. Pallinger}} 
&
\includegraphics[width=\linewidth]{./images/pallinger-christoph-aass/Anton_32986_v2-AASS-0112mm.jpg}
\Captionof{figure}{Sauerstoffmaske mit Reservoir. \SourcePicture{Ch. Pallinger}}
\end{tabularx}
\end{WidePar}
\end{table}


\MassnahmenNG{Spezielle Maßnahmen}{Sauerstoffberieselung}
{\label{m:sauerstoffberieselung}}
{%
\begin{enumerate}
  \item Kontraindikationen und Gegenanzeigen prüfen:
  \begin{itemize}
    \item COPD (\prettyref{topic:copd})
    \item Hyperventilationssyndrom, Hyperventilationstetanie
    (\ref{topic:hyperventailationssyndrom})
  \end{itemize}
  \item Situationsgerechte Dosierung je nach
  zugrundeliegender Erkrankung oder Symptomen
  \item Auswahl des Hilfsmittels:
  \begin{itemize}
    \item \E{Sauerstoff\-brille}
    \ce{O2}-Flow {\textless}~5\LPerMin*,
    \item \E{Sauerstoff\-maske}
    \ce{O2}-Flow ab 5\LPerMin*,
    \item \E{Sauerstoff\-maske mit
    Reservoir} \ce{O2}-Flow ab 5\LPerMin*,
  \end{itemize}
  \item Patientenaufklärung
  \item Sauerstoffgerät einschalten und Durchflußrate
  einstellen
  \item Bei Verwendung einer Sauerstoff\-maske mit
  Reservoir: Reservoir füllen
  \item Sauerstoffbrille/-maske positionieren 
\end{enumerate}
}


\subsubsection{Sauerstoffzufuhr bei der Beatmung}

\Layout{0}{\TxBeschreibung}{%
Bei der Beatmung kann Sauerstoff zugeführt werden. Beim
Beatmungbeutel schließt man dazu einen \ce{O2}-\EE{Verbindungsschlauch} an den
Druckminderer und den Beatmungsbeutel an. Wenn Vorhanden kann ein
\ce{O2}-\EE{Reservoir} verwendet werden, dadurch wird die
Sauerstoffkonzentration auf fast 100\PC* erhöht. 

Bei Beatmungsgeräten kann -- je nach Typ -- das Mischungsverhältnis von
Sauerstoff und Umgebungsluft eingestellt werden.\ZFootnote{Z.\,B. Schalthebel
"`Air Mix"' bei Weinmann Medumat{\texttrademark} Standard.}
}{%
\begin{itemize}
  \item Beamtungsbeutel:
  \begin{itemize}
    \item \ce{O2}-\EE{Verbindungsschlauch}
    \item \ce{O2}-\EE{Reservoir}
  \end{itemize}
  \item Beatmungsgerät: einstellbar
\end{itemize}
}{}{}{}















\clearpage % TODO temp clearpage
\section{Beatmung}
\label{topic:beatmung-hilfsmittel}
\index{Beatmung}
\index{Beatmung!Grundsätzliches}

\Layout{2}{\TxBeschreibung}
{Bei \EE{fehlender} oder \EE{nicht ausreichender
Spontanatmung} (Eigenatmung) wird eine \EE{künstliche
Beatmung} durchgeführt. Dabei wird Eigenatmung des Patienten
künstlich durch geeignete Hilfsmittel ersetzt oder 
ergänzt. Es handelt es sich um eine
\EE{Überdruckbeatmung}\ZFootnote{Theoretisch kann auch
mittels einer Unterdruckkammer eine Unterdruckbeatmung
durchgeführt werden (\I[noindex]{Eiserne Lunge}), dies ist
aber sehraufwendig und wird praktisch nicht durchgeführt.},
d.\,h. es wird mit Druck Luft in die Atemwege des Patienten
gepumpt. (Im Unterschied dazu wird bei der natürlichen
Atmung Luft mit Unterdruck in den Patienten eingesogen.)


Man Unterscheidet zwischen einer \E{kontrollierten} und
einer \E{assistierten} Beatmung. 

Bei der \T[Beatmung!kontrollierte]{kontrollierten Beatmung}
wird die gesamte Atmung künstlich aufrecht erhalten. 

Bei der
\T[Beatmung!assistierte]{assistierten Beatmung} hat der
Patient eine eigene Spontanatmung, die aber nicht ausreichend ist (zu langsam
oder zu seicht). Man unterstützt dabei die Atembemühungen
des Patienten, indem man zusätzliche oder tiefere Atemhübe
verbreicht. 

Eine Beatmung durch Fachpersonal sollte wenn
möglich immer mit Hilfsmitteln wie dem \E{Beatmungsbeutel} oder einem
\E{Beatmungsgerät} durchgeführt werden.
}{}{}{}{}

\Layout{0}{\TxGefahren}
{Der notwendige \EE{Überdruck} kann zu Problemen führen, wie
z.\,B. Überschreiten des \E{Speiseröhren-Öffnungsdrucks}
(\EE{Magenbeatmung}, \EE{Magenblähung}, Erbrechen) oder
Lungenschäden.

Es ist oft zu beobachten, das (wahrscheinlich in Folge von
Stress) Beatmungen zu schnell bzw. zu tief erfolgen. Dies
kann katastrophale Auswirkungen haben, da es zu einer
\EE{Hyperventilation} und damit zu einer atmungsbedingen
Alkalose kommt (Störung des Säure-Basen-Haushaltes, vgl.
\prettyref{topic:alkalose-respiratorische}).

\Cave{Cave: Hyperventilation und atmungsbedingen Alkalose}

\Fazit{Faustregel: Eine kontrollierte Beatmung eines
Erwachsenen soll grundsätzlich im Eigenrhythmus des Helfers erfolgen: Helfer
atmet \Sign{→} Patient wird beatmet.}
}{
\begin{itemize}
  \item Zu hoher Druck \Sign{→} Magenbeatmung/-blähung
  \item Zu schnell, zu tief: Hyperventialtion,
  atmungsbedingte Alkalose
\end{itemize}
}{}{}{}

\Layout{0}{Hilfmittel}
{Eine Beatmung wird normalerweise mit Hilfsmitteln
durchgeführt. Herlferseitig kommen dabei ein 
\EE{Beatmungsbeutel} (\prettyref{topic:beatmungsbeutel})
oder ein \EE{Beatmungsgerät} zum Einsatz. 
Patientenseitig werden \EE{Beatmungsmasken} oder
\EE{Beatmungsschläuche} (Endotrachealtuben\Z{,
Larynxtuben, -masken, Combituben, \ldots}) angewendet.

\bigskip

\bigskip
}{
\begin{itemize}
  \item Helferseitig:
  \begin{itemize}
    \item Beatmungsbeutel
    \item Beatmungsmaske
  \end{itemize}
  \item Patientenseitig
  \begin{itemize}
    \item Beatmungsmaske
    \item Beatmungsschlauch
  \end{itemize}
\end{itemize}
}{}{}{}


\subsection{Beatmungsbeutel}
\label{topic:beatmungsbeutel}
\index{Beatmungs!-beutel}



\Layout{2}{\TxBeschreibung}
{Der Beatmungsbeutel ist das einfachste Hilfsmittel,
welches dem Fachpersonal zur Verfügung steht, und kommt in
vielen Bereichen (Rettungsdienst, Krankenhaus, \ldots) zum
Einsatz. Er ermöglichst die kontrollierte oder
unterstützende Beatmung eines Patienten. Er kann mittels
einer Gesichtsmaske (Beatmungsmaske) angewendet, oder an
andere Syysteme (Tuben) angeschlossen werden. }
{%
\begin{itemize}
  \item Hilfsmittel zur kontrollierten und unterstützenden Beatmung
\end{itemize}
}
{}
{}
{}

\Layout{37}{Abbildung}
{}
{}
{./images/pallinger-christoph-aass/Ambubeutel_32741-AASS-0176mm.jpg}
{Beatmungsbeutel "`Ambu{\texttrademark} Mark IV"' mit
Reservoir, Bakterienfilter, PEEP-Ventil und Beatmungsmaske} {Ch.
Pallinger}

% \Layout{22}{Bestandteile}
% {
% Ein Beatmungsbeutel besteht im Idealfall aus folgenden Teilen:
% \begin{itemize}
% 	\item Beatmungsmaske 
% 	\item Bakterienfilter
% 	\item Ein-/Ausatemventil mit angeschlossenem PEEP-Ventil
% 	\item Beatmungsbeutel
% 	\item Reservoir, welches mit einem Verbindungsschlauch
% 	(\E{\ce{O2}-Line}) mit einer Berieselungseinheit verbunden
% 	ist.
% \end{itemize}
% }{}{}{}

\Layout{0}{Bestandteile}{%
Ein Beatmungsbeutel besteht üblicherweise aus einem
\EE{Luftzufuhrventil}, einem \EE{Balg} und einem
\EE{T-Stück} mit einem Ventil.

An der Luftzufuhr kann eine \EE{Sauerstoffleitung} und/oder
ein \EE{Sauerstoffreservoir} angeschlossen werden, um den
Sauerstoffgehalt des Einatemgases zu erhöhen. Am einfachsten
erfolgt dies über einen Verbindungsschlauch und eine
\ce{O2}-Berieselungseinheit, welche mit einem Flow von 15\LPerMin*{}
eingestellt ist. Mit Hilfe eines Reservoirs, welches
Sauerstoff während der Beatmungspausen zwischenspeichert,
kann die \ce{O2}-Konzentration des Einatemgases weiter
gesteigert werden.~\ZFootnote{Andere \ce{O2}-Systeme sind am
Markt verfügbar (Demand-Ventile, etc.), aber nur begrenzt
verbreitet.}

\Fazit{Wenn vorhanden soll der Beutel mit einer
Sauerstoffzufuhr verbunden werden.}

Am T-Stück gibt es \E{3 Schenkel}: Einer kommt vom
Beatmungsbalg, einer führt zum Patienten und einer zum
Auslass. Das \EE{Ventil} des T-Stückes \E{trennt die Ein- von
der Ausatemluft} des Patienten. Am Patientenschenkel kann ein
\EE{Bakterienfilter} und Messvorrichtungen (z.\,B. zur
\ce{CO2}-Messung) angebracht werden. Er wird mit einer
\EE{Beatmungsmaske} oder einem Tubus -- je nach Bedarf --
verbunden. Am Ausführungsschenkel kann ein \EE{PEEP}-Ventil
aufgesteckt werden.
}{%
\begin{itemize}
  \item Luftzufuhrventil
  \begin{itemize}
    \item \ce{O2}-Leitung
    \item Reservoir
  \end{itemize}
  \item Balg
  \item T-Stück
  \begin{enumerate}
    \item Vom Balg kommend
    \item Zum/vom Patienten: + Bakterienfilter, Beatmungsmaske od. Tubus
    \item Ausführung, evtl. mit PEEP-Ventil
  \end{enumerate}
\end{itemize}
}
{./images/geraete/ambu-komplett-zerlegt.jpg}
{Bestandteile eines
Beatmungsbeutels.}
{}

\begin{table}[H]
\Captionof{table}{Mögliche Sauerstoffkonzentrationen bei der Beatmung.}
\FontTableBase
\begin{tabulary}{\linewidth}{Lc}
%\addlinespace 
\noindent\TabHeading{Art der Beatmung} &
\noindent\TabHeading{\ce{O2}-Konzentration}
\\\addlinespace\midrule
Mund-zu-Mund Beatmung
&
17\PC*
\\\addlinespace
Beutel-Masken-Beatmung
&
21\PC*
\\\addlinespace
Beutel-Masken-Beatmung mit angeschlossenem \ce{O2} (15\LPerMin*)
&
\VonBis{40}{70}\PC*
\\\addlinespace
Beutel-Masken-Beatmung mit angeschlossenem \ce{O2} (15\LPerMin*) und
Reservoir &
fast 100\PC*
\\
\end{tabulary}
\hrule
\end{table}


\Layout{57}{}{}{}{
\begin{tabularx}{\linewidth}{XXX}
{\includegraphics[width=\linewidth]{./images/pallinger-christoph-aass/Ambubeutel_32741-AASS-0112mm.jpg}
\Captionof{figure}{Beamtungsbeutel "`Ambu{\texttrademark} Mark
IV"' mit Reservoir, \ce{O2}-Verbindungsschlauch,
Bakterienfilter und aufgestecktem PEEP-Ventil.
\SourcePicture{Ch. Pallinger}}} & {\includegraphics[width=\linewidth]{./images/pallinger-christoph-aass/Beatmungsmaske_32906_v2-AASS-0112mm.jpg}
\Captionof{figure}{Eine \E{Beatmungs}maske.
\SourcePicture{Ch. Pallinger}}} &
{\includegraphics[width=\linewidth]{./images/pallinger-christoph-aass/Mundstuecke_32723-AASS-0112mm.jpg}
\Captionof{figure}{Beatmungsmasken unterschiedlicher Größen:
0, 2, 3/4, 5. Die Farben der Masken haben
keine besondere Bedeutung. \SourcePicture{Ch. Pallinger}}}
\\\addlinespace\bottomrule
\end{tabularx}
}
{Bilderserie: Beatmungsbeutel und Beatmungsmasken.}
{}



\Layout{0}{Hygienische Wiederaufbereitung}
{%
Es gibt Beatmungsbeutel als Einweg- und Mehrwegprodukt.
Einwegprodukte müssen nach Gebrauch entsorgt werden.

Mehrwegprodukte können i.\,d.\,R. durch Desinfektion im
Einlegebad\ZFootnote{Autoklavierung ist meistens auch möglich
und wünschenswert, sofern die entsprechenden Einrichtungen
vorhanden sind.} wiederaufbereitet werden. Es ist unbedingt
notwendig, dass der Beatmungsbeutel dazu \EE{komplett
zerlegt} wird!

Die Herstellerhinweise sind zu beachten.
}{
\begin{itemize}
  \item Einweg- und Mehrwegprodukte
  \item I.\,d.\,R. Desinfektion mittels Einlegebad
  \begin{itemize}
    \item \Ca{Beutel komplett zerlegen!}
  \end{itemize}
\end{itemize}
\subsubsection{Technik der Beutel-Masken-Beatmung}

\TODO{EMHO}{2010-08-04}{Beatmungstechnik: Text fehlt}{}

\TakeNotesLarge

}{}{}



\subsubsection{Besonderheit: Beatmung während der
Reanimation}

Die Beatmung während der Reanimation ist eine kontrollierte
Beatmung. Wird eine Beatmungsmaske verwendet, so ist das
Verhältnis zur Herzdruckmassage zu beachten (vgl.
\CO{[AED, EHI]}; \prettyref{chp:AED}, \ref{chp:EHI}). Sobald
die Beatmung über einen Beamtungsschlauch (Tubus) erfolgt,
ist sie unabhängig von der Herzmassage.

\TakeNotesLarge












\subsection{Beatmungsgeräte}
\index{Beatmungs!-gerät}
\index{Beatmungs!-gerät!Medumat{\texttrademark}}
\index{Medumat{\texttrademark}}
\index{Medumat{\texttrademark}!Standard a}
\index{Medumat{\texttrademark}!Compact}
\index{Medumat{\texttrademark}!Standard}
\index{Medumat{\texttrademark}!Transport}

\begin{figure}[H]
\begin{center}
\includegraphics[width=7.189cm,height=4.94cm]{./images/beatmungsgeraet-medumat.png}
\Captionof{figure}{Beatmungsgerät \emph{Medumat
Standard\texttrademark}. \SourcePicture{ASB
Floridsdorf-Donaustadt}}
\end{center}
\end{figure}




Mit dem Beatmungsgerät \emph{"`Medumat"'} der Fa. Weinmann
kann ein Patient sowohl beatmet als auch berieselt werden
(sofern ein Berieselungsmodul vorhanden ist). 

Dafür gibt es zwei Bedienfelder: Das eine zum Berieseln,
welches vom Sanitäter eingestellt und verwendet werden
kann; das zweite darf nur auf Anweisung eines Notarztes
verwendet werden und wird für das Beatmen benötigt. 

\opt{ORGASBFLO}{
\begin{minipage}{\linewidth}
\ASBFLO{Beim ASB Floridsdorf-Donaustadt kommen folgende Geräte
zum Einsatz:

\begin{itemize}
  \item \EE{Notfallbeatmungsgeräte}:
  \begin{itemize}
    \item Hersteller: Fa. Weinmann, Produktfamilie:
    \T{Medumat}: 
    \begin{inparaitem}
\item Medumat{\texttrademark} Compact
\item Medumat{\texttrademark} Standard
\item Medumat{\texttrademark} Standard a
\end{inparaitem}
  \end{itemize}
  \item \EE{(Intensiv-)Transportbe\-atmungs\-ge\-rät}\nocite{ThierbachInterhospitaltransfer1}: \begin{itemize}
\item Hersteller: Fa.
Weinmann, Produktfamilie: \T{Medumat{\texttrademark}}:
\begin{inparaitem}
  \item Medumat{\texttrademark} Transport
\end{inparaitem}
\end{itemize}
\end{itemize}
}
\end{minipage}
}



\begin{figure}[H]
\includegraphics[width=\linewidth]{./images/geraete/ger-medumat-standard-bedienfeld.jpg}
\CaptionofDiff{figure}{Bedienfelder Berieselungseinheit \E{Modul Oxygen} (li.)
und Beatmungsgerät \E{Medumat{\texttrademark} Standard} (re.).

\E{Modul Oxygen (li.):} Flowmeter, Ein-Ausschalter und
Regelventil. Der Anschluss links dient der Verbindung der
Einheit mit dem Sauerstoffnetz des Fahrzeuges.

\E{Medumat{\texttrademark} Standard (re.):}
\Z{Air-Mix-Schalter, Ein-Ausschalter, Stellräder für
Atemfrequenz, Atemminutenvolumen (!) und Druckbegrenzung,
Beatmungsdruckanzeige, Warnleuchten für Verlegung
("`Stenosis"'), Lösung des Beatmungsschlauches
("`Disconnection"'), niederen Flaschendruck ("`<\,2.7\Bar*"')
und schwache Batterie.} } {Bedienfelder Berieselungseinheit
\E{Modul Oxygen} (li.) und Notfallbeatmungsgerät
\E{Medumat{\texttrademark} Standard} (re.).}
\end{figure}



\subsection{PEEP und PEEP-Ventil}
% \index{PEEP-Ventil}

\Definition{\T{PEEP} steht für "`Positive End-Exspiratory
Pressure"' (positiver end-exspiratorischer Druck), d.\,h.
positiver Druck in den Atemwegen am Ende der
Ausatmungsphase.}

\Layout{20}{\TxBeschreibung{}}
{%
Der \T{PEEP} sorgt dafür, dass am Ende der Ausatmungsphase
(Exspiration) der Druck in den Atemwegen nicht auf 0
absinkt, sondern stattdessen ein positiver Druck aufrecht
erhalten wird. Dadurch wird das Zusammenfallen der
Lungenbläschen (Alveolen, \prettyref{topic:alveolen})
verhindert.

\Z{Der Druck wird je nach Modell in \MetricUnit{cm}
Wassersäule (\MetricUnit{cm\,\ce{H2O}}) oder Millibar
(\MetricUnit{mbar}) angebeben. Typische Werte sind z.\,B. 5
oder 10\,\MetricUnit{cm\,\ce{H2O}}.}

Ein \Abk{PEEP} kann mittels eines
\T[PEEP-Ventil]{PEEP-Ventils} erreicht werden. Es wird dabei
auf die Ausatemöffnung des Beatmungsbeutels bzw.
Beatmungsschlauches aufgesteckt. Durch Drehen der Kappe kann
der gewünschte \Abk{PEEP} eingestellt werden. Die
Einstellung kann an der angebrachten \E{Skala} abgelesen
werden. Es gibt \Abk{PEEP}-Ventil-Produkte sowohl als
\EE{Einweg-}, als auch als \EE{Mehrweg}material! \Z{Moderne
Beatmungsgeräte verfügen oft über eine integrierte
PEEP-Funktion und benötigen kein zusätzliches
\Abk{PEEP-Ventil.}

Bei der Maskenbeatmung ist zu beachten, dass der PEEP nur
erreicht werden kann, wenn das Beatmungssystem \E{luftdicht}
ist. Sobald die Maske nicht mehr dicht aufsitzt, verliert
das \Abk{PEEP}-Ventil seine Wirkung, da die Ausatemluft
unkontrolliert über das \E{Leck} entweichen kann.

\Cave{Die Anwendung des \Abk{PEEP} ist grundsätzlich eine
\EE{ärztliche Maßnahme}, daher muss es \EE{bei der Anwendung
durch nicht-ärztliches Personal auf Null} gestellt sein. } 
}
}
{}{./images/pallinger-christoph-aass/Peep_32909-AASS-0112mm.jpg}
{PEEP-Ventil.}
{Ch. Pallinger}











\section{Hebe- und Transporttechniken}
\index{Hebetechniken}
\index{Transporttechniken}

\TODO{GABS}{2010-08-14}{Hebe- und Transporttechniken: Text
fehlt}{} 
\Layout{2}{{\TxQuerverweise}}{%
\begin{inparaitem}
  \item Tragsessel: \prettyref{topic:tragsessel};
  \item Fahrtrage, "`Einheitskrankentrage"': \ref{topic:fahrtrage};
\end{inparaitem}
}{}{}{}{}

Das Fachpersonal hat während des gesamten Transports die
Verantwortung für den Patienten, dies gilt auch für gehfähige
Patienten. Sollte der Patient auf Grund seines hohen Alters
oder seiner Kreislaufsituation nicht selbstständig gehen
können, muss er entweder \EE{sitzend} oder \EE{liegend}
transportiert werden. 

Patienten, bei denen mit einer Zustandsverschlechterung zu rechnen ist, sollen
grundsätzlich auf der Trage transportiert werden, da ein Umlagern im Fahrzeug
sehr umständlich und zeitraubend ist. 




\subsubsection{Umlagern von Patienten}
\index{Umlagern}
\label{topic:umlagern}

{\TakeNotesLarge}


\subsubsection{Umgang mit einem Rollstuhl}
\index{Rollstuhl}
\label{topic:rollstuhl}

% \prettyref{fig:rollstuhl-umgang}

{\TakeNotesLarge}

%     \begin{center}
% 
%     \includegraphics[width=10cm]{./images/noauth/GER-umlagern.png} 
% 
%     \Captionof{figure}{\AuthNotesPicLegalNotOk{GER-umlagern.png}{}\label{fig:rollstuhl-umgang}\label{fig:umlagern}Umlagern von Patienten.}
%     \end{center}



\subsubsection{Tragen mit dem Tragering}

\label{topic:tragering}

Das Tragen mit dem Tragering funktioniert nur bei bewusstseinsklaren und
kooperativen Patienten, außerdem sollte das Stiegenhaus, oder die zurück
zulegende Strecke, breit genug sein. Die Technik eignet sich nur für kurze
Distanzen.

\Layout{57}{}
{}{}
{
\begin{tabularx}{\linewidth}{XXX}
\includegraphics[width=\linewidth]{./images/motal-aass/00800/tragering1b.jpg}
&
\includegraphics[width=\linewidth]{./images/motal-aass/00800/tragering2b.jpg}
&
\includegraphics[width=\linewidth]{./images/motal-aass/00800/tragering3b.jpg}
\\
Tragering &
Transport mit dem Tragering &
Transport mit dem Tragering 
\\\addlinespace\bottomrule
\end{tabularx}
}{Bilderserie: Transport mit einem Tragering}{}




\subsubsection{Tragetuch}

\index{Tragetuch}
\label{topic:bergetuch}
\label{topic:tragetuch-taetigkeit}

Das Tragetuch (ESE-Tuch) eignet sich ideal für den Transport
von liegenden Patienten wo der Einsatz einer Krankentrage
eventuell aus Platzgründen nicht möglich ist. Es wird auch
gerne zum Umlagern von Patienten oder zum Retten aus einer
Gefahrenzone verwendet. Die Füße des Patienten zeigen
talwärts, in der Ebene wird der Patient wird immer mit den
Füßen in Transportrichtung getragen.

Der Transport mit Hilfe des Tragetuchs ist jedoch nicht
wirbelsäulenschonend. Ist eine Immobilisation der
Wirbelsäule notwendig soll stattdessen eine andere Technik
(Vakuummatratze, etc.) angewandt werden.)

\Layout{57}{}
{}{}
{
\begin{tabularx}{\linewidth}{XXX}
\includegraphics[width=\linewidth]{./images/motal-aass/00800/tragetuch1b.jpg}
&
\includegraphics[width=\linewidth]{./images/motal-aass/00800/tragetuch2b.jpg}
&
\includegraphics[width=\linewidth]{./images/motal-aass/00800/tragetuch3b.jpg}
\\
Tuch aufbreiten \ldots &
\ldots die näher zum Patient liegende Hälfte
umschlagen\ldots 
& \ldots die obere Hälfte nochmals
umschlagen\ldots 
\\\addlinespace
\includegraphics[width=\linewidth]{./images/motal-aass/00800/tragetuch4b.jpg}
&
\includegraphics[width=\linewidth]{./images/motal-aass/00800/tragetuch5b.jpg}
&
\includegraphics[width=\linewidth]{./images/motal-aass/00800/tragetuch6b.jpg}
\\
\ldots und das Tuch an den Patienten legen. 
&
Patient leicht anheben und Tuch bis zur Hälfte
unterschieben\ldots 
& \ldots Patient schonend zurück legen.
\\\addlinespace
\includegraphics[width=\linewidth]{./images/motal-aass/00800/tragetuch7b.jpg}
&
\includegraphics[width=\linewidth]{./images/motal-aass/00800/tragetuch8b.jpg}
&
\includegraphics[width=\linewidth]{./images/motal-aass/00800/tragetuch9b.jpg}
\\
Von der anderen Seite ebenfalls leicht anheben und
Tuch ausbreiten. 
& 
Zu dritt das Tuch fest anpacken\ldots 
&
\ldots und auf Kommando gemeinsam anheben
\\\addlinespace\bottomrule
\end{tabularx}
}{Bilderserie: Transport mit dem Tragetuch\index{Tragetuch!Bilderserie}}{}

% \begin{figure}[H]
%     \begin{center}
%  TODO BILD: Bilderserie: Tragetuch 
%     \includegraphics[width=10cm]{./images/noauth/GER-tragetuch-tragen.png}
% 
%     \Captionof{figure}{\AuthNotesPicLegalNotOk{GER-tragetuch-tragen.png}{}}
%     \end{center}
% \end{figure}

\clearpage % TODO temp clearpage
\subsection{Lagerung}
\subsubsection{Lagerungsarten}
\label{topic:lagerungsarten}
\index{Lagerungsarten}

\Layout{57}{}
{}{}
{
\begin{tabularx}{\linewidth}{XXX}
\includegraphics[width=\linewidth]{./images/motal-aass/00800/lagerungflachb.jpg}
&
\includegraphics[width=\linewidth]{./images/motal-aass/00800/lagerung30gradb.jpg}
&
\includegraphics[width=\linewidth]{./images/motal-aass/00800/lagerungoberkoerperhochb.jpg}
\\
Flachlagerung &
Lagerung mit 30\textdegree{} erhöhtem Oberkörper &
Lagerung mit stark erhöhtem Oberkörper 
\\\addlinespace
\includegraphics[width=\linewidth]{./images/motal-aass/00800/lagerungseiteb.jpg}
&
\includegraphics[width=\linewidth]{./images/motal-aass/00800/lagerungkollaps2b.jpg}
&
\includegraphics[width=\linewidth]{./images/motal-aass/00800/lagerungkollaps1b.jpg}
\\
Seitenlagerung & 
Lagerung mit erhöhten Beinen. 

\Z{Diese Lagerung wurde früher
als in Hinblick auf den Hypovolämischen Schock als
\I{Schocklagerung} bezeichnet. Da sich die Lagerung je nach Schockart unterscheidet (Die Lagerung mit erhöhten
Beinen ist beim kardiogenen Schock lebensgefährlich!), soll diese Bezeichnung nicht
mehr verwendet werden.}
& Lagerung mit
erhöhten Beinen bei Verdacht auf Wirbelsäulenverletzung
\\\addlinespace
\includegraphics[width=\linewidth]{./images/motal-aass/00800/lagerungabdomenb.jpg}
&
\includegraphics[width=\linewidth]{./images/motal-aass/00800/lagerungthromboseb.jpg}
&
\includegraphics[width=\linewidth]{./images/motal-aass/00800/lagerungembolieb.jpg}
\\
Lagerung zur Entspannung der Bauchdecke &
Lagerung mit erhöhter Exxtremität &
Lagerung mit herabhängender Gliedmaße
\\\addlinespace\bottomrule
\end{tabularx}
}{Bilderserie: Lagerungsarten}{}

\subsubsection{Dekubitusprophylaxe}
\index{Dekubitus!Prophylaxe}
\label{topic:dekubitusprophylaxe}
\TODO{EMHO}{2010-08-04}{Dekubitusprophylaxe: Text fehlt}{}

\TakeNotesLarge



\clearpage
\section{Hygienische Händedesinfektion}
\label{topic:hygienische-haendedesinfektion}
% \marginpar{\includegraphics[width=4cm]{images/AASdev20090228-img102.jpg}
% 
% \includegraphics[width=4cm]{images/AASdev20090228-img101.jpg}}
\Definition{Desinfektion der Hände und Handgelenke mittels
einer Desinfektionslösung und einer genormten
Handwaschtechnik.}

% \Married{
\begin{WidePar}

\begin{multicols}{2}
\paragraph{Warum?} 
Die Hände des Personals sind die
häufigsten Keimüberträger. 

Die normale Händewaschung entfernt nur groben Schmutz und
sonstige mechanische Verunreinigungen. Die
Desinfektionslösung reduziert gezielt aktiv viele Keime. 

Versuche mit Fingerfarben
zeigen, dass mit einer herkömmlichen Handwaschtechnik große
Teile der Hand ausgelassen werden (Fingerzwischenräume,
Handgelenk, Daumenkuppe,\ldots).

\begin{minipage}{\linewidth}

\paragraph{Ziel}
\begin{itemize}
    \item Keimreduktion
    \item Vermeidung von \EE{Kreuzinfektionen}
    \begin{itemize}
        \item Patient1 ${\rightarrow}$ Sanitäter ${\rightarrow}$ Patient 2
        \item \EE{"'Unterbrechung der
        Infektionskette{\textquotedblright}}
    \end{itemize}
\end{itemize}
\paragraph{Wann durchzuführen?}
\begin{itemize}
    \item Dienstantritt
    \item \EE{Nach jedem Patienten!!}
\end{itemize}
\paragraph{Womit?}~
\begin{itemize}
    \item Händedesinfektionslösung:
    Sterilium\texttrademark{}, Manopronto\texttrademark{}
    o.\,ä.
\end{itemize}
\end{minipage}
\end{multicols}
\end{WidePar}
% } {
% \begin{center}
% \includegraphics[width=3.961cm,height=3.169cm]{images/AASdev20090228-img102.jpg}
% 
% \includegraphics[width=4.001cm]{images/AASdev20090228-img101.jpg}
% \end{center}
% }

\begin{WidePar}

\begin{minipage}{\linewidth}

\paragraph{Vorgehen}\vspace{1em}~ 

\begin{center}
\includegraphics[width=0.23\linewidth]{./images/hygiene/haendedesinfektion01.png}\hspace{0.01\linewidth}
\includegraphics[width=0.23\linewidth]{./images/hygiene/haendedesinfektion02.png}\hspace{0.01\linewidth}
\includegraphics[width=0.23\linewidth]{./images/hygiene/haendedesinfektion03.png}\hspace{0.01\linewidth}
\includegraphics[width=0.23\linewidth]{./images/hygiene/haendedesinfektion04.png}

\includegraphics[width=0.23\linewidth]{./images/hygiene/haendedesinfektion05.png}\hspace{0.01\linewidth}
\includegraphics[width=0.23\linewidth]{./images/hygiene/haendedesinfektion06.png}\hspace{0.01\linewidth}
\includegraphics[width=0.23\linewidth]{./images/hygiene/haendedesinfektion07.png}\hspace{0.01\linewidth}
\includegraphics[width=0.23\linewidth]{./images/hygiene/haendedesinfektion08.png}
\end{center}
\Captionof{figure}{Hygienische Händedesinfektion. \SourcePicture{DEMAW}}
\vspace{1em}
\begin{multicols}{2}

\begin{enumerate}
    \item Bei Dienstantritt: Ringe, Uhren, Armbänder,
    Schmuckstücke, etc. entfernen, Handwaschung
    durchführen, um grobe Verunreinigungen zu entfernen.
    \item Eine Portion alkoholisches
    Händedesinfektionsmittel (ca. 3\Ml* = 1~Hohlhand) aus
    dem Spen\-der entnehmen.
    \item Mittels Ellenbogen, vollkommene Benetzung beider Hände und
    Handgelenke.
    \item Handinnenflächen gegeneinander reiben
    \item Geschlossene Finger.
    \item Handgelenke einreiben
    \item Mit ineinander verschränkten Fingern Handinnenflächen
    gegeneinander reiben.
    \item Mit rechter Handinnenfläche linken Handrücken reiben und
    umgekehrt, dabei Finger ineinander verschränken.
    \item Hände ineinander verhaken und Finger gegeneinander bewegen
    \item Daumen mit gegenüberliegender Hand vollständig umschließen
    und rotierend reiben.
    \item Daumenkuppen nicht vergessen!
    \item Fingerkuppen im Handteller kreisförmig einreiben bis Alkohol
    verdunstet ist
\end{enumerate}
\EE{Einwirkzeit beachten!} (lt. Herstellervorgaben, i.d.R. mind. 30 Sekunden) 
\end{multicols}
\end{minipage}
\end{WidePar}

\nocite{KirschnickPflege2}






\section{Manuelle Traumatechniken}
% Michael Motal

\Layout{2}{{\TxQuerverweise}}{%
\begin{itemize}
  \item Traumacheck: \prettyref{topic:traumacheck}, Seite
  \pageref{topic:traumacheck}.
\end{itemize}
}{}{}{}{}

\subsection{Helmabnahme und manuelle HWS-Stabilisierung}
\label{topic:helmabnahme}
\label{topic:hws-stabilisierung-manuell}
\index{Helmabnahme}
\index{HWS!manuelle Stabilisierung}


\Layout{2}{\TxBeschreibung}
{Ein Patient mit aufgesetztem Sturzhelm stellt eine eine
besondere Herausforderung dar. Der Sturzhelm ist
grundsätzlich immer abzunehmen! }{}{}{}{}


\Layout{2}{Material und Personal}{
\begin{itemize}
\item 2 Helfer ({\color{DefaultA}\E{hinterer} Helfer},
{\color{DefaultB}\E{seitlicher} Helfer})
\item 1 HWS-Schiene für später
\end{itemize}
}{}{}{}{}


% \begin{multicols}{2}
\Layout{2}{Vorgehen}{
\begin{enumerate}
  \item \SpeechMargin{"`Bitte bleiben Sie
  ruhig liegen und bewegen Sie sich nicht!"'}\EE{Dem Patient
  in Blickrichtung nähern, um
  unnötige Bewegungen des Kopfes zu vermeiden, und
  gleichzeitig auffordern den Kopf nicht zu bewegen.}
    \item {\color{DefaultA}\EE{Fixieren}: Der \E{hintere}
    Helfer fixiert den Helm von hinten und lässt nicht mehr los, bis der Helm entfernt ist. Zwischen Helm und Schritt eine Helmhöhe
    Abstand lassen, damit der Helm in einer Bewegung
    abgezogen werden kann}
    \item {\color{DefaultB}\EE{Ansprechen}: Der
    \E{seitliche} Helfer spricht den Patienten von vorne an, öffnet
    Visier und Kinnriemen}
    \item {\color{DefaultB}Der \E{seitliche} Helfer
    \EE{stabilisiert} den Kopf während der Helmabnahme: eine
    Hand fasst mit Daumen und Zeigefinger das Unterkiefer,
    die andere Hand stützt den Kopf von unten.}
    \item {\color{DefaultA}\EE{Helm abziehen}: Der
    \E{hintere} Helfer zieht den Helm in einer
    \EE{S-förmigen Bewegung} ab: den Helm \EE{zuerst kippen
    bis die Nasenspitze} sichtbar ist, dann die Unterkante
    in der \EE{gegengleichen Bewegung} zum Helfer bewegen.}
    
    {\color{DefaultB}Während der Helm abgezogen wird,
    rutscht der \E{seitliche} Helfer mit deiner Haltehand
    etwas nach.}
    \item {\color{DefaultA}Der \E{hintere} Helfer zieht den
    Helm sanft fertig ab und legt ihn sicher ab.}
    \item {\color{DefaultA}\EE{Der \E{hintere} Helfer
    übernimmt den Kopf} im doppelten C-Griff: Daumen und
    Zeigefinger um ein Ohr. Der Kopf darf
    \EE{\E{nicht} auf Zug} gehalten werden, die \Abk{HWS}
    darf \EE{nicht gestaucht} werden.
    Der Kopf wird in
    Neutralposition
    stabilisiert. \cite{ITLSde6,AsboeLehrmeinung2011-07}
    \Commentary[HWS-Immobilisation / Zug]{ASBÖ: Nach der Lehrmeinung 07/2011 vom 01.02.2011 \cite{AsboeLehrmeinung2011-07} darf der Kopf nicht mehr unter Zug gehalten werden. Dies ist konsistent mit \Abk{ITLS} (6.\,A.)
    \cite{ITLSde6}}}
    \item {\color{DefaultB}Der \E{seitliche} Helfer legt
    eine \Abk{HWS}-Schiene an
    (\FullPrettyRef{topic:hws-schiene-anlegen}).}
\end{enumerate}
% \end{multicols}
}{}{}{}{}

\begin{figure}[H]
\FormatTable
\begin{WidePar}
\Captionof{figure}{Bilderserie: Helmabhname und Anlage
einer HWS-Schiene}

\subfloat
[Während des Annäherns den Patient auffordern, den Kopf nicht zu bewegen. ]
{\includegraphics[width=0.32\linewidth]{./images/motal-aass/00800/helm1.jpg}}
\hfill
\subfloat[Den Kopf sofort mit den Händen ficieren, Visier öffnen und Patienten ansprechen.]
{\includegraphics[width=0.32\linewidth]{./images/motal-aass/00800/helm2.jpg}}
\hfill
\subfloat[Ggfs. Patient synchron und schonend auf den Rücken drehen.]
{\includegraphics[width=0.32\linewidth]{./images/motal-aass/00800/helm3.jpg}}


\subfloat[Ggfs. Brille entfernen, Kinnriemen öffnen und behindernde
Kleidung entfernen. ]
{\includegraphics[width=0.32\linewidth]{./images/motal-aass/00800/helm4.jpg}}
\hfill
\subfloat[Halswirbelsäule stabilisieren.]
{\includegraphics[width=0.32\linewidth]{./images/motal-aass/00800/helm5.jpg}}
\hfill
\subfloat[Helm kippen bis von oben die Nasenspitze sichtbar ist,
danach vorsichtig nach hinten wegziehen ]
{\includegraphics[width=0.32\linewidth]{./images/motal-aass/00800/helm6.jpg}}

\subfloat[Nach dem Abziehen Stabilisierung des Kopfes übernehmen. Keinen Zug und keine Stauchung ausüben!]
{\includegraphics[width=0.32\linewidth]{./images/motal-aass/00800/helm7.jpg}}
\hfill
\subfloat
[HWS-Schiene anlegen]
{\includegraphics[width=0.32\linewidth]{./images/motal-aass/00800/helm8.jpg}}
\hfill
\parbox{0.32\linewidth}{~}

\end{WidePar}
\end{figure}

% \Layout{57}{}
% {}{}
% {
% \begin{tabularx}{\linewidth}{XXX}
% \includegraphics[width=\linewidth]{./images/motal-aass/00800/helm1.jpg}
% &
% \includegraphics[width=\linewidth]{./images/motal-aass/00800/helm2.jpg}
% &
% \includegraphics[width=\linewidth]{./images/motal-aass/00800/helm3.jpg}
% \\
% Während des Annäherns den Patient auffordern, den Kopf
% nicht zu bewegen. 
% &
% Helm sofort fixieren,Visier öffnen und Patient ansprechen.
% & 
% Ggfs. Patient synchron und schonend auf den Rücken drehen.
% \\\addlinespace
% \includegraphics[width=\linewidth]{./images/motal-aass/00800/helm4.jpg}
% &
% \includegraphics[width=\linewidth]{./images/motal-aass/00800/helm5.jpg}
% &
% \includegraphics[width=\linewidth]{./images/motal-aass/00800/helm6.jpg}
% \\
% Ggfs. Brille entfernen, Kinnriemen öffnen und behindernde
% Kleidung entfernen. 
% &
% Halswirbelsäule stabilisieren.
% & 
% Helm kippen bis von oben die Nasenspitze sichtbar ist,
% danach vorsichtig nach hinten wegziehen 
% \\\addlinespace
% \includegraphics[width=\linewidth]{./images/motal-aass/00800/helm7.jpg}
% &
% \includegraphics[width=\linewidth]{./images/motal-aass/00800/helm8.jpg}
% &
% \\
% Nach dem Abziehen Stabilisierung des Kopfes übernehmen.
% Keinen Zug und kene Stauchung ausüben!
% & 
% HWS-Schiene anlegen
% (\FullPrettyRef{topic:hws-schiene-anlegen} 
% &
% \\\addlinespace\bottomrule
% \end{tabularx}
% }{Bilderserie: Helmabnahme}{}



\clearpage % TODO temp clearpage
\section{Schienung}
\label{topic:schienung}

\subsection{Anlegen einer
HWS-Immobilisationsschiene (Stifneck\texttrademark)}
\label{topic:hws-schiene-anlegen}\index{HWS-Schiene!anlegen}\index{Stifneck|see{HWS-Schiene}}
\label{topic:hws-schiene-beschreibung}
\index{HWS-Schiene!Beschreibung}

\begin{compactdesc}
  \item[Häufige Markennamen:] Stifneck{\texttrademark},
select{\texttrademark}Stifneck select{\texttrademark},
Stifneck Paedi select{\texttrademark}
\end{compactdesc}


\Layout{60}{\TxBeschreibung}
{%
Dient zum Ruhigstellen und Stabilisieren der {\bfseries H}als{\bfseries
W}irbel{\bfseries S}äule (HWS). Sie soll bei jedem Verdacht auf eine
Verletzung der Halswirbelsäule eingesetzt werden; auf jeden
Fall nach einer Helmabnahme und bei Verdacht auf ein Schädel-Hirn-Trauma
(\Abk{SHT})!

Es sind verschiedene
Fabrikate im Einsatz (z.\,B. Stifneck\texttrademark). Es gibt
Schienen mit einer fixen Größe oder moderne Schienen, bei
welchen sich die gewünschte Größe einstellen lässt
(Stifneck select{\texttrademark} bzw. Stifneck Paedi
select{\texttrademark} für Kinder).

}{

}
{./images/noauth/GER-stiffneck.jpg}
{\AuthNotesPicLegalNotOk{GER-stiffneck.jpg}{}HWS-Schiene.}
{}



\Layout{2}{Material und Personal}{
\begin{itemize}
    \item 2 Helfer ({\color{DefaultA}\E{hinterer} Helfer},
{\color{DefaultB}\E{seitlicher} Helfer})
    \item 1 \Abk{HWS}-Schiene (z.\,B. Stifneck\texttrademark)
\end{itemize}
}{}{./images/noauth/GER-stiffneck.jpg}{}{}


% \begin{multicols}{2}
\Layout{2}{Vorgehen}{
\begin{enumerate}
    \item \EE{Patient auffordern den Kopf nicht zu bewegen.}
    \item {\color{DefaultA}Der hintere
    Helfer \EE{stabilisiert} den Kopf und die \Abk{HWS}
    in Neutralposition ohne Zug und ohne Stauchung.}
    \item {\color{DefaultB}Beengende und behindernde
    \EE{Kleidung entfernen}.}
    \item {\color{DefaultB}\EE{Größenmessung}: Der seitliche
    Helfer misst die Entfernung von Schulter bis Unterkiefer in Fingerbreiten
    aus.}
    \item {\color{DefaultB}Stifneck entsprechend einstellen:
    gemessene Entfernung einstellen.
    
    Verschiebemechanismus des
    Stifneck blockieren, ab nun darf kein Verstellen mehr
    möglich  sein}
    \item {\color{DefaultB}\Abk{HWS}-Schiene \EE{zuerst am
    Kinn anlegen} und festhalten}
    \item {\color{DefaultB}Klettverschlussband so einfalten,
    dass es nicht am Boden schleifen kann (hält sonst nicht
    mehr)}
    \item {\color{DefaultB}Nackenteil unter Patient durch
    führen (Vorsicht: Klettverschlussband)}
    \item {\color{DefaultB}Fest schließen.}
    \begin{itemize}
        \item[] Kinnspitze sollte ein kurzes Stück über Kante hinaus schauen
        (kein Durchrutschen) aber nicht allzu weit (falsch eingestellt)
    \end{itemize}
\end{enumerate}
% \end{multicols}
}{}{}{}{}

\Layout{57}{}
{}{}
{
\begin{tabularx}{\linewidth}{XXX}
\includegraphics[width=\linewidth]{./images/motal-aass/00800/hwsschiene1.jpg}
&
\includegraphics[width=\linewidth]{./images/motal-aass/00800/hwsschiene2.jpg}
&
\includegraphics[width=\linewidth]{./images/motal-aass/00800/hwsschiene3.jpg}
\\
Nach der Helmabnahme folgt sofort das Anlegen einer
HWS-Schiene. 
&
Der Kopf bleibt in Neutralposition zug- und
stauchungsfrei fixiert. D

ie passende Größe der HWS-Schiene
wird bestimmt und eingestellt. 
& 
HWS-Schiene zuerst am Kinn anlegen und danach den
Nackenteil unter dem Patient durch führen.
\\\addlinespace\bottomrule
\end{tabularx}
}{Bilderserie: Anlegen einer HWS-Schiene}{}

\clearpage % TODO temp clearpage
\subsection{Schaufeltrage}
\label{topic:schaufeltrage-anwendung}
\label{topic:schaufeltrage-handhabung}
\index{Schaufeltrage}

\index{Schaufeltrage!Beschreibung}
\label{topic:schaufeltrage-beschreibung}

\begin{compactdesc}
  \item[Häufige Markennamen:] Ferno ScoopEXL{\texttrademark}
  \item[Anwendung:] \prettyref{topic:schaufeltrage-anwendung}
\end{compactdesc}

\Layout{60}{\TxBeschreibung}
{%
Die Schaufeltrage wird primär für das Umlagern von liegenden
Patienten mit dem Verdacht auf eine Verletzung der
Wirbelsäule, Becken oder Oberschenkel verwendet. Sie
ermöglich eine möglichst bewegungsfreie und Schmerzarme
Rettung bzw. Umlagerung. Das Anwendungsgebiet ist jedoch
sehr vielfältig. Es reicht vom sicheren Transport im
unwegsamen Gelände bis zum Retten aus einer LKW-Kabine. 

Letztendlich soll
der Patient zwecks Schienung auf einer \EE{Vakuummatratze}
gelagert werden.
Je nach Fabrikat kann der Patient jedoch auch direkt auf
einer Schaufeltrage (ohne
Verwendeung einer Vakuummatratze) immobilisiert werden, dazu
werden spezielle Gurte und Kopfstützen benötigt. 

Die Schaufeltrage setzt sich aus 2 Längshälften zusammen;
die Fußteile sind längenverstellbar. Die Kopfteile haben eine fixe Länge.
Sollte der Patient größer
als die Trage sein, müssen Kopf und Oberkörper auf jeden
Fall trotzdem auf der Trage zu liegen kommen. Ausnahme: Beinverletzung, bei
der eine Wirbelsäulenverletzung ausgeschlossen werden kann.

\Z{Das Spineboard
(\FullPrettyRef{topic:spineboard-beschreibung}) ist in
vielen Fällen eine Alternative zur Schaufeltrage
\cite{ITLSde6,PHTLS6,AsboeLehrmeinung2011-10}.}
}{

}
{./images/noauth/GER-schaufeltrage.jpg}
{\AuthNotesPicLegalNotOk{GER-schaufeltrage.jpg}{}Eine Schaufeltrage}
{}

\subsubsection{Aufschaufeln mittels Schaufeltrage}
\label{topic:schaufeltrage-aufschaufeln}

\Layout{2}{Material und Personal}{
\begin{itemize}
    \item Mind. 2 Helfer
    \item Schaufeltrage
    \item Gurte für Schaufeltrage
\end{itemize}
}{}{}{}{}


%%%%%%%%%%%%%%%%%%%%%%%%%%%%%%%%%%%%%%%%%%%%%%%%%%%%%%%%%%%
% TODO : Vervollständigen
\TODO{GABS}{2010-mm-dd}{Schaufeltrage, Kopfdrehung: Vervollständigung
notwendig.}
{Vervollständigung notwendig! \par Dieser Teil ist unvollständig, 
   er benötigt eine Überarbeitung!    
   Idealerweise sollte ein 3ter Sanitäter dabei sein um den Kopf mitzudrehen.
   Das Mitdrehen des Kopfes kommt aber gar nicht vor bzw. wird nicht
   beschrieben. }
%%%%%%%%%%%%%%%%%%%%%%%%%%%%%%%%%%%%%%%%%%%%%%%%%%%%%%%%%%%


% \begin{multicols}{2}
\Layout{2}{Vorgehen}{
\begin{enumerate}
    \item Schaufeltrage neben Patienten in zusammengebautem Zustand auf die
    richtige Länge einstellen. Sollte ein Patient zu groß sein
    müssen Kopf und Oberkörper auf jeden Fall trotzdem auf der Trage
    liegen! Vorsicht: Markierung der max. Länge bei Schraubmodell
    \item \EE{Öffnen} der Trage: Schnappverschlüsse; je nach Methode nur den
    unteren oder beide öffnen.
    \item \EE{Aufschaufeln}:
    \begin{enumerate}
      \item 1. Methode: Fußseite offen, Patient wird von oben
    kommend wie auf eine Schere aufgeladen, die Trage unter ihm
    geschlossen. Kontrolle des sicheren Schlusses!
    
    Diese Methode ist besonders bei Becken- und Oberschenkelverletzungen
    sinnvoll und der zweiten Methode vorzuziehen da der Patient nicht am Becken
    bewegt wird.
    \item 2. Methode: beide Teile getrennt, Patient leicht an
    Schulter und Hüfte zur Seite gedreht und jeweiligen Tragenteil
    darunter geschoben. Hälften zuerst am Kopf-, dann am Fußende
    schließen. Kontrolle des sicheren Schlusses!
    \end{enumerate}
    \item Angurten: mehrere Gurte\footnote{Anzahl der Gurte lt.
    Herstellerangabe. Wir empfehlen die Verwendung von mind. \E{3} Gurte};
    jeden Gurt zuerst um den Handgriff der einen, dann der anderen Hälfte
    führen,  dann über dem Patienten schließen und festziehen.
    \item Umlagern auf Vakuummatratze: Trage mit vorbereiteter Matratze
    (\prettyref{topic:vakuummatratze-anwendung}) neben Patienten stellen um diesen
    einen möglichst geringen Weg tragen/bewegen zu müssen.
    \item Gurte öffnen und entfernen.
    \item Hälften trennen: beide Schnappverschlüsse öffnen
    \item Hälften parallel zum Patienten flach wegziehen. Der Patient soll
    hierbei nicht bewegt werden!
    \item Sollten die Schnappschlösser klemmen kann man die Trage wie bei
    der Scherentechnik an nur einer Seite öffnen und in V-Form (gleichzeitig, je
    ein Helfer auf einer Seite) entfernen.
\end{enumerate}
% \end{multicols}
}{}{}{}{}

\subsubsection{Umlagern mittels Schaufeltrage}
\label{topic:schaufeltrage-umlagern}

\Layout{2}{Material und Personal}{
\begin{itemize}
    \item Mind. 2 Helfer
    \item Patient auf Schaufeltrage
    (\FullPrettyRef{topic:schaufeltrage-aufschaufeln})
\end{itemize}
}{}{}{}{}


%%%%%%%%%%%%%%%%%%%%%%%%%%%%%%%%%%%%%%%%%%%%%%%%%%%%%%%%%%%
% TODO : Vervollständigen
\TODO{GABS}{2010-mm-dd}{Schaufeltrage, Kopfdrehung: Vervollständigung
notwendig.}
{Vervollständigung notwendig! \par Dieser Teil ist unvollständig, 
   er benötigt eine Überarbeitung!    
   Idealerweise sollte ein 3ter Sanitäter dabei sein um den Kopf mitzudrehen.
   Das Mitdrehen des Kopfes kommt aber gar nicht vor bzw. wird nicht
   beschrieben. }
%%%%%%%%%%%%%%%%%%%%%%%%%%%%%%%%%%%%%%%%%%%%%%%%%%%%%%%%%%%


% \begin{multicols}{2}
\Layout{2}{Vorgehen}{
\begin{enumerate}
  \item Patient liegt, wie unter
  \FullPrettyRef{topic:schaufeltrage-aufschaufeln},
  beschrieben, auf der Schaufeltrage.
    \item Umlagern, z.\,B. auf Vakuummatratze: Trage mit
    vorbereiteter Matratze (\prettyref{topic:vakuummatratze-anwendung}) neben Patienten stellen um diesen
    einen möglichst geringen Weg tragen/bewegen zu müssen.
    \item Gurte öffnen und entfernen.
    \item Hälften trennen: beide Schnappverschlüsse öffnen
    \item Hälften parallel zum Patienten flach wegziehen. Der Patient soll
    hierbei nicht bewegt werden!
    \item Sollten die Schnappschlösser klemmen kann man die Trage wie bei
    der Scherentechnik an nur einer Seite öffnen und in V-Form (gleichzeitig, je
    ein Helfer auf einer Seite) entfernen.
\end{enumerate}
% \end{multicols}
}{}{}{}{}


\subsubsection{Immobilisation auf der Schaufeltrage}
\label{topic:schaufeltrage-immobilisation}

\Commentary[Schaufeltrage / Immobilisation]{Die
Immobilisation direkt auf der Schaufeltrage ist -- entgegen
vieler Gerüchte -- nicht nur möglich, sondern auch von
Herstellern, wie zum Beispiel Ferno{\texttrademark}
ausdrücklich freigegeben (z.\,B. Ferno
\E{ScoopEXL{\texttrademark}} Schaufeltrage). Die Produkte
sind auch entsprechnd konstruiert (Röntgendurchlässigkeit, \ldots).
Zubehör, wie z.\,B. die Ferno
\E{SpiderStraps{\texttrademark}} oder Ferno \E{Head
Immobilizer for 65EXL{\texttrademark}}, werden ausdrücklich
für Schaufeltragen angeboten.

Ein entsprechendes Instruktionsvideo kann unter
\url{http://www.ferno.com/scoopexl/scoopexl.htm} abgerufen
werden.  }

\ASBFLO{Die Immobilisation auf der Schaufeltrage wird nicht
empfohlen, da die aufnehmenden Spitäler i.\,d.\,R. nicht
darauf vorbereitet sind (Tauschmaterial, \ldots)}

\Layout{2}{Material und Personal}{
\begin{itemize}
    \item Schaufeltrage mit Zulassung zur Immobilisation
    von Patienten
    \item Passende Kopfstützen
    \item Passendes Gurtsystem zur Immobilisation (z.\,B.
    Ferno \E{Fastrap{\texttrademark} Quick Restraint
    System}, \FullPrettyRef{topic:spiderstraps-beschreibung})
    \item Patient auf Schaufeltrage
    (\FullPrettyRef{topic:schaufeltrage-aufschaufeln})
\end{itemize}
}{}{}{}{}




% \begin{multicols}{2}
\Layout{2}{Vorgehen}{
\begin{enumerate}
  \item Patient liegt, wie unter
  \FullPrettyRef{topic:schaufeltrage-aufschaufeln},
  beschrieben, auf der Schaufeltrage.
    \item Gurtsystem und Kopfstützen wie vom Hersteller
    beschrieben anbringen. 
    
    Je nach Schaufeltrage und Gurtsystem kann die Technik
    erheblich abweichen.
\end{enumerate}
% \end{multicols}

\Cave{Bevor der Kopf fixiert wird, muss der Körper bereits
verläßlich fixiert sein!} 
}{}{}{}{}

%%%%%%%%%%%%%%%%%%%%%%%%%%%%%%%%%%%%%%%%%%%%%%%%%%%%%%%%%%%
% TODO Sandwich-Technik: Neuerstellung
\TODO{GABS}{yyyy-mm-dd}{Sandwich-Technik: Neuerstellung von Grafik/Text 
                          notwendig}{}
%%%%%%%%%%%%%%%%%%%%%%%%%%%%%%%%%%%%%%%%%%%%%%%%%%%%%%%%%%%

\clearpage % TODO temp clearpage
\subsection{Vakuummatratze}
\label{topic:vakuummatratze-anwendung}\index{Vakuummatratze!Handhabung}
\label{topic:vakuummatratze-beschreibung}





\Layout{2}{\TxBeschreibung}{
Allgemeine
Beschreibung: \prettyref{topic:vakuummatratze-beschreibung}

Die Vakuummatratze dient zur Ruhigstellung des gesamten Körpers bei Verdacht
auf mögliche Wirbelsäulenverletzungen, Becken- oder Oberschenkelfrakturen
(Oberschenkelhals und Schaft).
Am Ende soll die Matratze an die Körperform der Patienten eng
angepasst sein und somit Bewegungen effektiv verhindern.

\ASBFLO{Die Immobilisation auf der Vakuummatratze ist die
Methode der Wahl für die Wirbelsäulenimmobilisation.
\Commentary[Vakuummatratze / Immobilisation]{Die Immobilisation auf der Vakuummatratze
wird derzeit als Methode der Wahl empfohlen, da die Technik
organisationsübergreifend weit verbreitet und das
Personal entsprechend geschult ist. Aufnehmende Spitäler
sind auf Patienten, welche auf einer Vakuummatratze
immobilisiert wurden, gut vorbereitet. }

}}{}{}{}{}

\Layout{27}{Abbildung}
{}
{}
{./images/noauth/GER-vakuummatratze.png}
{Vakuummatratze}
{}

\Layout{2}{Material und Personal}{
\begin{itemize}
    \item Mind. 2 Helfer
    \item Trage
    \item Vakuummatratze
    \item Absauggerät
    \item Leintuch
    \item Evtl. Adapter, je nach Matratzentyp
\end{itemize}
}{}{}{}{}


% \begin{multicols}{2}
\Layout{2}{Vorgehen}{
\begin{enumerate}
    \item Fahrtrage auf tiefste Stufe stellen (ganz unten).
    \item Vorbereiten: Vakuummatratze auf Fahrtrage legen. Die weichere
    Seite ist die Patientenseite, Ventil beim Kopf. Kugeln so verteilen,
    dass der gesamte Körper gut, sowie die betroffene Stelle besonders gut
    fixiert werden kann.
    \item Vorabsaugen mit Absauggerät bis die Matratze gut formbar aber nicht allzu
    fest ist.
    \item Vorformen, Ventil verschließen. Wenn vorhanden Leintuch auf
    die Matratze legen.
    \item Angegurteten Patienten mit Schaufeltrage auf die Vakuummatratze
    legen, Schaufeltrage entfernen (\prettyref{topic:schaufeltrage-handhabung}).
    \item \EE{Absaugen} mit Absauggerät, dabei zügig Matratze an den Patienten
    \EE{anpassen}. Besonderes Augenmerk gilt der Fixierung der verletzten
    Regionen sowie dem Kopf ("`Ohren"'
    formen: Ecken neben dem Kopf hoch formen). Einige Modelle bieten die Möglichkeit einer zusätzlichen Fixierung mittels an der Vakuummatratze befestigter Gurte.
    
    \Cave{\Ca{Vorsicht!} Der Kopf und die HWS dürfen beim Anpassen nicht gestaucht werden!}
    \item Ist kein weiteres Absaugen mehr möglich und die Matratze gut
    angepasst (Matratze muss jetzt bretthart sein) Ventil schließen,
    Absauggerät abdrehen.
    \item Patient und Vakuummatratze an Fahrtrage angurten.
\end{enumerate}
% \end{multicols}
}{}{}{}{}

\subsection{Spineboard}
\label{topic:spineboard-anwendung}\index{Spineboard!Anwendung|Z}
\label{topic:spineboard-beschreibung}\index{Spineboard!Beschreibung|Z}

\Layout{2}{\TxBeschreibung}{
Allgemeine Beschreibung:
\prettyref{topic:spineboard-beschreibung}
Das Spineboard ist im Prinzip eine flaches Brett auf dem
der Patient ähnlich dem Schaufeltrage-Vakuummatratze-Duett
gerettet und fixiert werden kann. Es stellt somit in vielen
Situationen eine Alternative zur
Ganzkörper-/Wirbelsäulenimmobilisation dar.

Diese Technik wurde
ursprünglich vor allem im anglo-amerikanischen Raum
eingesetzt und erfährt auch hierzulande immer größere
Verbreitung.

\ASBFLO{Das Spineboard gehört derzeit nicht zur
Standardausstattung und wird nicht routinemäßig verwendet.}
}{}{}{}{}

\ASBFLO{Das Spineboard gehört derzeit nicht zur
Standardausstattung und wird nicht routinemäßig verwendet.}




% \TakeNotesMedium

\subsection{Gurtensysteme für Schaufeltrage und Spineboards}

\subsubsection{Ferno Fastrap{\texttrademark} Quick Restraint System}
\index{SpiderStraps!Beschreibung}
\label{topic:spiderstraps-beschreibung}
\index{Fastrap{\texttrademark}!Beschreibung}
\label{topic:fastrap-beschreibung}

\Layout{2}{\TxBeschreibung}
{%
% Anwendung: \prettyref{topic:schaufeltrage-anwendung}
Das Fastrap{\texttrademark}-System kann sowohl zusammen mit
einer Schaufeltrage, als auch mit einem Spineboard verwendet
werden. \Commentary[Ferno Fastrap{\texttrademark} /
Schaufeltrage]{\Speech{"`The Fastrap system is compatible
with Ferno backboards and Scoop stretchers."'} Ferno International Online Shop,
\url{http://www.fernointernational.com}, 2011-05-13}
% 
% {\TakeNotesLarge}
}{

}
{}
{}
{}

\Layout{27}{Abbildung}
{}
{}
{./images/pallinger-christoph-aass/Spiderstrap_32954_v2-AASS-0112mm.jpg}
{Ferno Fastrap{\texttrademark} Quick Restraint System}
{Ch. Pallinger}










% 
\subsection{Rettungskorsett (KED{\texttrademark}-System)}
\label{topic:rettungskorsett-anwendung}\index{Rettungskorsett!Anwendung}\index{KED-System!anlegen}
\index{Rettungskorsett!Beschreibung}\label{topic:rettungskorsett-beschreibung}

\begin{compactdesc}
  \item[Häufige Markennamen:] Kendrick Extrication
  Device{\texttrademark} (KED{\texttrademark})
  \item[Anwendung:] \prettyref{topic:rettungskorsett-anwendung}
\end{compactdesc}



\Layout{2}{\TxBeschreibung}
{%
Das Rettungskorsett wird primär für das Umlagern von
sitzenden Patienten mit dem Verdacht auf eine Verletzung der
Wirbelsäule verwendet. Es kann aber auch für das Retten
aller sitzenden Patienten aus einer Gefahrenzone verwendet
werden. Am Kopfgriff des
Korsetts kann, falls nötig, ein Kranhaken eingehängt
werden. Das Rettungskorsett wird (ausser wenn es
ausschliesslich dazu verwendet wird, den Patienten mit
"`Haltegriffen"' auszustatten, gemeinsam mit einer
\EE{HWS-Schiene} verwendet.
}{

}
{}
{}
{}

\Layout{27}{Abbildung}
{}
{}
{./images/geraete/KED.jpg}
{Ein Rettungskorsett.}
{Wikipedia, "`Onthost"', PD}


\Layout{2}{Material und Personal}{
\begin{itemize}
    \item 2\,--\,3 Helfer
    \item Rettungskorsett
    \item HWS-Schiene
    \item Vorbereitete Vakuummatratze
\end{itemize}
}{}
{}
{}
{}


% \begin{multicols}{2}
\Layout{2}{Vorgehen}{
\begin{enumerate}
    \item Patient aufsuchen und sofort HWS-Schiene anlegen.
    (\prettyref{topic:hws-schiene-anlegen})
    \item Ein Helfer stabilisiert von nun an am besten von hinten den  Kopf um
    seitliche Bewegungen des Kopfes zu vermeiden.
    \item Rettungskorsett \E{vorsichtig} in aufgefaltenem Zustand
    hinter dem Patienten platzieren. Die Oberkanten der
    Flügel sollen in der Achsel des Patienten zu liegen
    kommen.
    \item Die Beingurte vom Korsett lösen und baumeln
    lassen.
    \item Der mittlere und untere Bauchgurt wird
    geschlossen und halbfest angezogen. Das Korsett soll in der Höhe nicht mehr
    verrutschen können.
    \item Kopf mit der zum System gehörigen faltbaren
    Schaumstoffunterlage so unterpolstern, dass die HWS in
    Neutralposition steht.
    \item Der Kopf wird mit den Klettbändern fixiert: ein
    Band an der Stirn, das andere unter dem Kinn die HWS-Schiene
    entlang. Die Klettbänder werden am Rettungskorsett
    fixiert. Der Kopf soll nun seitlich nicht mehr geneigt
    werden können. Der Patient wird \E{unter keinen
    Umständen zu einem Versuch aufgefordert!}
    \item Die Beingurte werden geschlossen: von unten um
    den Oberschenkel herum nach oben; vorsichtig hin- und
    zurückziehend so eng es geht anlegen und festzurren.
    Vorsicht auf den Genitalbereich!
    \item Die halbfest geschlossenen Bauchgurte nun
    festzurren.
    \item Den letzten, obersten Bauchgurt schließen und
    festzurren. Dieser wird erst jetzt geschlossen um die
    Atemarbeit des Patienten nicht zu behindern.
    \item Der Patient ist nun entlang der gesamten
    Wirbelsäule geschient und kann mittels der 3 am Korsett
    angebrachten Griffe gehoben werden.
    \item Sanftes Umlagern auf die vorbereitete
    Vakuummatratze (\prettyref{topic:vakuummatratze-anwendung}).
    \item Beingurte leicht lockern.
    \item Den obersten Bauchgurt öffnen.
    \item Vakuummatratze anformen und sichern (\prettyref{topic:vakuummatratze-anwendung}).
       
\end{enumerate}
% \end{multicols}
}{}{}{}{}


\subsection{Vakuum-Beinschiene}
\label{topic:vacuumbeinschiene-anwendung}
\index{Vakuumbeinschiene!Anwendung}
\label{topic:vacuumbeinschiene-beschreibung}

\Layout{2}{\TxBeschreibung}
{%
Die Vakuumbeinschiene dient zum Ruhigstellen von Verletzungen
\EE{unterhalb des Knies}. Verletzungen des Oberschenkels
können mit ihr \E{nicht} ruhiggestellt werden, da sie das Hüftgelenk nicht schient
(\Sign{→}\,Vakuummatratze)!
% 
% {\TakeNotesLarge}
}{

}
{}
{}
{}

\Layout{37}{Abbildung}
{}
{}
{./images/pallinger-christoph-aass/Beinschiene_32977_v2-AASS-0176mm.jpg}
{Vakuumbeinschiene mit Zubehör}
{Ch. Pallinger}


\Layout{2}{Material und Personal}{
\begin{itemize}
    \item 2 Helfer: Einer seitlich, einer beim Fußende
    \item Vakuumbeinschiene
    \item Absauggerät
    \item Evtl. Adapter
\end{itemize}
}{}{}{}{}


% \begin{multicols}{2}
\Layout{2}{Vorgehen}{
\begin{enumerate}
    \item Länge anpassen, gegebenenfalls oben umfalten.
    Klettverschlussbänder herrichten.
    \item Kugeln gut verteilen. Auch am Fußende müssen genug Kugeln
    sein, um dem Fuß sowohl unter der Sohle als auch am Fußrücken
    gut Halt geben zu können
    \item Schiene leicht vorabsaugen, damit Kugeln nicht mehr verrutschen
    können
    \item Schuh am gesunden Bein ausziehen: Finden der schonensten Methode.
    Ein Helfer hält das Bein auf Zug, der andere öffnet den Schuh
    (Bänder ausfädeln, zur Not aufschneiden)
    \item vorsichtiges Wiederholen zu zweit am verletztem Bein, Socken
    ausziehen (DMS)
    \item Stiefelgriff: gleich nach Ausziehen von Schuh/Socken nimmt der
    Helfer am Fußende durch das Loch in der Schiene die Ferse des
    verletzten Fußes, die andere Hand hält den Fußrücken (Stiefelgriff).
    \EE{Das Bein wird so auf dem Boden liegend auf Zug gehalten, der Helfer
    verbleibt so, bis die Schiene fest sitzt.}
    \item Der seitliche Helfer richtet die Schiene vorsichtig unter dem
    Patienten aus und fixiert die Schiene mit den Klettverschlussbändern.
    Um den Fuß werden Ohren gebogen und kreuzweise festgeklettet, alternativ ein
    Band über den Fußrücken und eines unter der Sohle. Der Fuß soll von allen
    Seiten so gut wie möglich umschlossen sein!
    \item Der seitliche Helfer saugt ab und zieht dabei die Klettbänder
    immer wieder fest nach. \EE{Der andere Helfer hält nach wie vor den
    Fuß im Stiefelgriff.}
    \item Fertig abgesaugt: Ventil schließen, Absauggerät abschalten. Der
    Stiefelgriff kann jetzt losgelassen werden.
\end{enumerate}
% \end{multicols}
}{}{}{}{}

Die Schiene soll bretthart sein, das Bein leicht auf Zug fixiert. Der
Fuß soll weder seitlich noch zum Körper hin/vom Körper weg beweglich sein.



\subsection{Aluminiumkernschiene}
\label{topic:sam-splint-anwendung}
\label{topic:schienung-arm}
\index{Sam-Splint{\protect\texttrademark}!Anwendung}
\index{Aluminiumkernschiene!Anwendung}
\label{topic:sam-splint-beschreibung}
\index{Aluminiumkernschiene!Beschreibung}
\index{Sam-Splint{\protect\texttrademark}!Beschreibung}

\begin{compactdesc}
  \item[Häufige Markennamen:]
  SAM-Splint{\protect\texttrademark}
\end{compactdesc}



\Layout{61}{}
{%
Der Aluminiumkernschiene ist eine flexible, vom Kunststoff
umgebene Metallschiene und ermöglicht das Ruhigstellen von
Extremitätenverletzungen mit abnormer Stellung. Sie lässt
sich leicht an den Arm anformen, muss jedoch noch mit einer
Mullbinde oder mit Dreiecktüchern fixiert werden. Die
\EE{Anpassung} erfolgt an der \EE{gesunden} Extremität.
}{

}
{./images/noauth/GER-samsplint.png}
{\AuthNotesPicLegalNotOk{GER-samsplint.png}{}Sam-Splint{\protect\texttrademark}.}
{}




\Layout{2}{Material und Personal}{
\begin{itemize}
    \item 2 Helfer
    \item Sam-Splint
    \item Peha-Haft, Mullbinde + Leukoplast oder
    \VonBis{2}{3} Dreiecktuchkrawatten
    \item Ein weiteres Dreiecktuch als Armtragetuch
\end{itemize}
}{}{}{}{}


\Layout{2}{Vorgehen}{
\begin{enumerate}
    \item Sam-Splint in der Hälfte falten, an gesundem Arm der Länge
    nach anpassen, Knick beim Ellbogen
    \item Anlegen am verletzten Arm, sanft anpassen. Fixierung mit
    Peha-Haft, Mullbinden oder Dreiecktuchkrawatten
    \item Fingerspitzen müssen sichtbar sein (DMS), gegebenenfalls Enden
    umfalten.
    \item Lagerung des geschienten Armes mittels Armtragetuch
\end{enumerate}
}{}{}{}{}



\section{Verbände}

\subsection{Druckverband}
\label{topic:druckverband}\index{Druckverband!anlegen}

\Layout{2}{\TxBeschreibung}{%
Der Druckverband dient der Stillung einer starken, d.\,h.
potentiell lebensbedrohlichen, Blutung.
}{}{}{}{}


\Layout{2}{Material und Personal}{%
\begin{itemize}
    \item \VonBis{1}{2} Helfer
    \item sterile Wundauflage je nach Wundgröße
    \item Druckkörper: z.\,B. Mullbinde
    \item Dreiecktuchkrawatte
\end{itemize}
}{}{}{}{}


% \begin{multicols}{2}
\Layout{2}{Vorgehen}{
\begin{enumerate}
    \item Sofort bei Erkennen einer \E{starken} Blutung durch
    direkten Druck auf die Wunde und Hochhalten die Blutung verlangsamen (Pat. kann dies je
    nach Zustand selbst tun!)
    \item Sterile Wundauflage auf die Wunde legen
    \item Druckkörper dem Wundverlauf folgend über die sterile
    Wundauflage legen, weiterhin drücken.
    \item Druckkörper mit einer Dreiecktuchkrawatte fest und möglichst
    vollständig umschließen.
    \item Dreiecktuchkrawatte festziehen und mit einem festen Knoten bei
    ausreichendem Druck direkt über der Wunde fixieren.
    \item 2. Knoten zur Sicherung des Ersten.
\end{enumerate}
% \end{multicols}
}{}{}{}{}

Sollte die Blutung zu einem späteren Zeitpunkt wieder
stärker werden oder der Druckverband zu locker angelegt
sein wird ein \EE{weiterer} Druckkörper mit einer
weiteren Dreiecktuchkrawatte noch fester angelegt. \EE{Ein
schon angelegten Druckverband wird nicht wieder geöffnet!}



% \subsection{Bilderserie}
% 
% \clearpage
% 
% \begin{WidePar}
% \begin{tabularx}{\linewidth}{XXXX}
% \toprule
% \includegraphics[width=\linewidth]{./images/khd/gef-wind-rauch.jpg}
% 
% Text 1
% &
% \includegraphics[width=\linewidth]{./images/khd/gef-wind-rauch.jpg}
% 
% Text 2
% &
% \includegraphics[width=\linewidth]{./images/khd/gef-wind-rauch.jpg}
% 
% Text 3
% &
% \includegraphics[width=\linewidth]{./images/khd/gef-wind-rauch.jpg}
% 
% Text 4
% \\\midrule
% \includegraphics[width=\linewidth]{./images/khd/gef-wind-rauch.jpg}
% 
% Text 1
% &
% \includegraphics[width=\linewidth]{./images/khd/gef-wind-rauch.jpg}
% 
% Text 2
% &
% \includegraphics[width=\linewidth]{./images/khd/gef-wind-rauch.jpg}
% 
% Text 3
% &
% \includegraphics[width=\linewidth]{./images/khd/gef-wind-rauch.jpg}
% 
% Text 4
% \\\midrule
% \includegraphics[width=\linewidth]{./images/khd/gef-wind-rauch.jpg}
% 
% Text 1
% &
% \includegraphics[width=\linewidth]{./images/khd/gef-wind-rauch.jpg}
% 
% Text 2
% &
% \includegraphics[width=\linewidth]{./images/khd/gef-wind-rauch.jpg}
% 
% Text 3
% &
% \includegraphics[width=\linewidth]{./images/khd/gef-wind-rauch.jpg}
% 
% Text 4
% \\\midrule
% \includegraphics[width=\linewidth]{./images/khd/gef-wind-rauch.jpg}
% 
% Text 1
% &
% \includegraphics[width=\linewidth]{./images/khd/gef-wind-rauch.jpg}
% 
% Text 2
% &
% \includegraphics[width=\linewidth]{./images/khd/gef-wind-rauch.jpg}
% 
% Text 3
% &
% \includegraphics[width=\linewidth]{./images/khd/gef-wind-rauch.jpg}
% 
% Text 4
% \\\midrule
% \includegraphics[width=\linewidth]{./images/khd/gef-wind-rauch.jpg}
% 
% Text 1
% &
% \includegraphics[width=\linewidth]{./images/khd/gef-wind-rauch.jpg}
% 
% Text 2
% &
% \includegraphics[width=\linewidth]{./images/khd/gef-wind-rauch.jpg}
% 
% Text 3
% &
% \includegraphics[width=\linewidth]{./images/khd/gef-wind-rauch.jpg}
% 
% Text 4
% \\\midrule
% \includegraphics[width=\linewidth]{./images/khd/gef-wind-rauch.jpg}
% 
% Text 1
% &
% \includegraphics[width=\linewidth]{./images/khd/gef-wind-rauch.jpg}
% 
% Text 2
% &
% \includegraphics[width=\linewidth]{./images/khd/gef-wind-rauch.jpg}
% 
% Text 3
% &
% \includegraphics[width=\linewidth]{./images/khd/gef-wind-rauch.jpg}
% 
% Text 4
% \\\midrule
% \\\bottomrule
% \end{tabularx}
% \end{WidePar}

\clearpage
\section{Assistenztätigkeiten}

\Layout{61}{}{%
\noindent Teile des Textes und der Handlungsanleitungen entstammen dem Skriptum
\emph{"`Ärztliche
Grundfertigkeiten"'} des \emph{Department für medizinische
Aus- und Weiterbildung -- Medizinische Universität Wien, Abt. Methodik und
Entwicklung},  welches in Zusammenarbeit  mit dem Klinischem Institut für
Hygiene  und Medizinische Mikrobiologie entstanden ist.

Wir danken Hrn. ao.Univ.Prof. Dr. Michael Schmidts und
seinem Team für die Unterstützung!
}{

}
{icons/muw.png}
{}
{MUW}









% \clearpage
\subsection{Aufziehen eines Medikamentes aus einer Ampulle}

\Layout{2}{\TxBeschreibung}{
Ampulle mit rotem Punkt: Sollbruchstelle, Daumen muss beim Aufbrechen
auf dem roten Punkt sein. 

Ampulle mit "`Halsring"': Ampulle kann
nach allen Richtungen aufgebrochen werden. 

% \marginpar{\scriptsize\H{Warum mach{\textquotesingle} ich mir wegen
% der Verletzungsgefahr Sorgen? Ist ja nur ein kleiner
% Schnitt, ist ja meine Sache {\dots}}
% 
% \H{\color{ubuntured}Nein! Unter Normalbedingungen, wenn ich in aller Ruhe mein
% Medi aufziehen kann, ist es tatsächlich mein eigenes Bier. Was aber wenn
% ich mich bei einer Reanimation an der Ampulle schneide? Die anderen
% Kollegen sind beschäftigt (Drücken, Blasen, Defi {\dots}) und dann
% steh ich da mit meinem kleinen Schnitt und blute alles an. Und der
% Notarzt braucht dringend das Adrenalin {\dots} Was jetzt?}
% 
% \Ca{Inakzeptabel!}}

Beim Arbeiten muss besonderes Augenmerk auf eine sterile Arbeitsweise gelegt
werden! Es kann sonst zu Infektionen kommen.
}{}{}{}{}


\Layout{2}{Material}{
\begin{compactitem}
  \item Spritze.
  \item Kanüle(n).
  \item Ampulle mit Medikament.
  \item Trockene Tupfer.
\end{compactitem}
}{}{}{}{}


% \begin{multicols}{2}
% \Layout{2}{Vorgehen}{
\begin{WidePar}
\MarginBoxFilledRounded{Vorgehen}{}
\bigskip
\begin{multicols}{2}
\begin{enumerate}
  \item Arzt verlangt ein Medikament.
  \item Sanitäter sucht das Medikament und zeigt die noch verschlossene Ampulle
  dem Arzt bevor sie geöffnet wird.
  \item \Z{Kontrolle Ablaufdatum}
  \item Kontrolle Inhalt: Verfärbung, Ausflockung?
  \item Öffnen der Ampulle:
  \begin{itemize}
    \item  Glasampulle:
    \begin{itemize}
      \item Ampullenspitze beklopfen, darauf achten, dass sich keine
      Flüssigkeit im Ampullenkopf mehr befindet
      \item Öffnen der Medikamenten-Ampulle mit Tupfer, roter Punkt weist
      zum Körper
      \marginpar{\includegraphics[width=\linewidth]{./images/ampulle_beklopfen.png}
      
      \Captionof{figure}{Ampullenkopf
      beklopfen. \SourcePicture{Bild: DEMAW}}}
    \end{itemize}
    \item Plastikampulle,
    "`Dreh-und-Drink"'-Verschluss:
    \begin{itemize}
      \item Den Verschlussknopf ohne Kontamination des Ampullenhalses
      (Sollbruchstelle) abdrehen.
    \end{itemize}
  \end{itemize}
  \item Ampulle abstellen.
  \item Verpackung der Spritze aufschälen, Spritze in Verpackung lassen.
  \item Verpackung der Kanüle aufschälen, Spritze auf Konus aufstecken.
  \item Entfernen der Kanülenabdeckung (nicht abdrehen, sondern
  abziehen!)
  \item Einführen der Kanüle in die Ampulle.
  \begin{itemize}
    \item Beim Einführen der Kanüle in die Ampulle den Ampullenrand
    nicht berühren.
  \end{itemize}
  \item Aufziehen des
  Medikamentes.\marginpar{\includegraphics[width=\linewidth]{./images/motal-aass/medi-aufziehen_4815-00800.jpg}}
  \item Vorsichtig entlüften.
  \item Variante "`Verabreichung über Venflon"':
  \begin{itemize}
    \item Nadel händisch abziehen und in den Kontamed verwerfen.
    \item Arzt das Medikament übergeben, Ampulle zeigen.
    \item Wird das Medikament nicht sofort verwendet, Spritze mit
    Stoppel verschließen und Ampulle ankleben.
  \end{itemize}
  \item Variante "`zuspritzen in eine Infusionsflasche"':
  \begin{itemize}
    \item Gummimembran desinfizieren.
    \item Kanüle nicht bis zum Schaft einstechen.
    \item Medikament zuspritzen.
    \item Beschriften der Flasche.
  \end{itemize}
  \item Variante "`direktes Spritzen"':
  \begin{itemize}
    \item Nach Kanülengröße fragen.
    \item Aufziehkanüle händisch entfernen.
    \item Spritzkanüle aufstecken.
    \item Spritze mit Verschlusskappe reichen, Ampulle zeigen.
  \end{itemize}
\end{enumerate}
\end{multicols}
\end{WidePar}
% }{}{}{}{}

\begin{table}[H]
\FormatTable
\begin{WidePar}
\Captionof{table}{Bilderserie: Medikament aus Ampulle
aufziehen}
\begin{tabularx}{\linewidth}{XXXX}
\includegraphics[width=\linewidth]{./images/motal-aass/medi-kontrolle_4807-00800.jpg}
\Captionof{figure}{Kontrolle \SourcePicture{Motal}}
&
\includegraphics[width=\linewidth]{./images/motal-aass/medi-aufziehen_4815-00800.jpg}
\Captionof{figure}{Aufziehen mit Nadel \SourcePicture{Motal}} 
&
\includegraphics[width=\linewidth]{./images/motal-aass/medi-kontrollieren_4824-00800.jpg}
\Captionof{figure}{Übergabe \SourcePicture{Motal}}
&
\includegraphics[width=\linewidth]{./images/motal-aass/medi-verabreichen_4826_CLOSEUP-00800.jpg}
\Captionof{figure}{Verabreichung über Venflon durch befugtes Personal \SourcePicture{Motal}}
\\\bottomrule
\end{tabularx}
\end{WidePar}
\end{table}

\subsection{Vorbereitung einer Infusion}
\label{topic:infusion}


\Layout{2}{Material}{%
\Married{
\begin{compactitem}
    \item Infusionsflasche.
    \item Infusionsbesteck.
    \item Aufhängevorrichtung.
    \item Alkoholhaltiges Desinfektionsmittel.
    \item einige saubere Tupfer.
    \item Permanentmarker.
\end{compactitem}
}{
% \centering\includegraphics[width=0.5\linewidth]{./images/infusion_tropfkammer.jpg}
}
}{}{}{}{}


\Layout{2}{Vorgehen}{%
Beim Arbeiten muss besonderes Augenmerk auf eine sterile Arbeitsweise gelegt
werden! Es kann sonst zu Infektionen kommen.

\begin{WidePar}
\begin{multicols}{2}
\begin{enumerate}
  \item Aufhängevorrichtung auf Flasche aufsetzen wenn nötig.
  \item Eröffnen der Flaschenabdeckung (Membranschutz).
  \item Wischdesinfektion der Gummimembran.
  \begin{itemize}
    \item Alkoholgetränkter Tupfer kann während weiterer Vorbereitung
    auf Membran belassen werden.
  \end{itemize}
  \item Montage des Infusionsbestecks
  \begin{itemize}
    \item \EE{Durchflußregler (Radklemme) schließen}
    \item Einstich des Tropfkammerdorns in Infusionsflasche mit
    geschlossenem Lüftungsventil und geschlossener Radklemme.
    \item "`No touch technique"' -- Dorn der
    Tropfkammer und Gummimembran der Infusionsflasche dürfen nicht
    berührt werden.
  \end{itemize}
  \item Infusionsschlauch füllen.
  \begin{itemize}
    \item Flasche wenden und hochhalten.
    \item Tropfkammerluftventil öffnen.
    \item Tropfkammer durch Zusammendrücken bis zur Hälfte füllen.
    \item Konusabdeckung des Infusionsschlauchs am anderen Ende muss bei der
    Belüftung nicht entfernt werden.
    \item Infusionsflasche und Ende des Infusionsschlauchs mit nicht
    dominanter Hand halten  .
    \item Radklemme öffnen.
    \item Infusionsschlauch entlüften.
    \item Radklemme schließen.
  \end{itemize}
\end{enumerate}
\end{multicols}
\end{WidePar}
}{}{}{}{}

\clearpage % TODO temp clearpage
\subsection{Venenverweilkanüle}\label{topic:venflon}

\Layout{2}{\TxBeschreibung}{
("`Venflon"'): Vorbereitung und Assistenz 
}{}{}{}{}

\begin{figure}[H]
\includegraphics[width=\linewidth]{./images/venflon.png}

\Captionof{figure}{Eine periphere Venenverweilkanüle
(Venflon\texttrademark), in ihre Einzelbestandteile zerlegt.
\SourcePicture{DEMAW}}
\end{figure}
\begin{table}[H]\FormatTable
\Captionof{table}{Verschiedene Größen von peripheren
Venenverweilkanülen. \Source{DEMAW (modifiziert)}}

\begin{tabular}{m{1.5619999cm}cccll}
\multicolumn{2}{l}{Nur zur Illustration!} 
&
Länge
&
 Durchfluss 
 &
 Farbcode
 &
\\
 {\O} mm & G & mm  & [\nicefrac{ml}{min}] & &
\\\midrule
\Z{0,5 x 0,8} & \Z{22} & \Z{25} & \Z{25} & blau & Sehr
dünn
\\
\Z{0,7 x 1,0} & \Z{20} & \Z{32} & \Z{55} & rosa & dünn
\\
\Z{0,9 x 1,2} & \Z{18} & \Z{40} & \Z{90} & grün &
normal
\\
\Z{1,1 x 1,4} & \Z{17} & \Z{42} & \Z{135} & weiß &
dick
\\
\Z{1,3 x 1,7} & \Z{16} & \Z{42} & \Z{170} & grau & Sehr
dick
\\
\Z{1,6 x 2,1} & \Z{14} & \Z{42} & \Z{265} & orange &
Pferd \\\bottomrule
\end{tabular}
\end{table}


% \begin{multicols}{2}
\Layout{2}{Material}{
\begin{compactitem}
  \item Nierentasse.
  \item Staubinde.
  \item Venflon:
  \begin{compactitem}
    \item Die Größe (= Farbe) ist vom Durchführenden vorher zu erfragen!
    \item Die Flügel werden hinunter geklappt.
    \item Die farbige Lasche des Zuspritzventils in rechten Winkel zur
    Venflonachse stellen.
  \end{compactitem}
  \item Einmalverschlusskappe ("`roter
  Stoppel"').
  \item Fixierungspflaster.
  \item einige saubere trockene Tupfer.
  \item Alko-Tupfer.
  \item (Einwegspritze mit 5\Ml* Spülflüssigkeit).
  \item (Faserschreiber).
  \item Kontamed-Box.
\end{compactitem}
% \end{multicols}
}{}{}{}{}


% \begin{multicols}{2}
\Layout{2}{Assistenz}{
\begin{itemize}
    \item Stauschlauch zureichen.
    \item Geöffnete Verpackung des Alko-Tupfers hinreichen, sodass dieser
    entnommen werden kann.
    \item Kontamed in Reichweite des  Durchführenden.
    bereitstellen, sodass er später die Nadel abwerfen kann.
    \item Venflon zureichen.
\begin{itemize}
    \item Venflon an der Schutzkappe halten, sodass der
    Durchführende den Venflon abziehen kann.
\end{itemize}
    \item Trockene Tupfer zureichen.
    \item Je nach Situation Infusionsschlauch, Stoppel oder gar nichts
    zureichen.
    \item Fixierpflaster zureichen.
\end{itemize}
% \end{multicols}
}{}{}{}{}

\A{Butterfly gehört erklärt}
% TODO Butterfly


\subsection{Intubation: Vorbereitung und Assistenz}
\label{topic:intubation-assistenz}


\Layout{22}{Material}
{
\begin{compactenum}
    \item Beatmungsbeutel inkl. passender Beatmungsmaske, Bakterienfilter und
    Sauerstoffreservoir.
    \item \ce{O2}-Berieselungseinheit und \ce{O2}-Line.
    \item Evtl. Beatmungsgerät.
    \item Absaugeinheit: \par\Ca{Keine Intubation ohne
    Absaugung!}
    \item Laryngoskop, bestehend aus:
    \begin{compactenum}
        \item Spatel: gebogen oder gerade; Größen 0\,--\,4.
        \item Griff.
    \end{compactenum}
    \item (Endotracheal-)Tubus; verschiedene Größen.
    \item Führungsdraht (\T{Mandrin}).
    \item Silikonspray.
    \item Evtl. Magillzange.
    \item Mit Luft gefüllte Spritze zum Cuffen (10\Ml*).
    \item Beißschutz: Beißkeil oder Guedel-Tubus.
    \item Fixationsmaterial (Mullbinde).
    \item Stethoskop.
\end{compactenum}
}
{}
{./images/pallinger-christoph-aass/Testlunge_32902-AASS-0112mm.jpg}
{Laryngoskopgriff und verschiedene Spatel}
{Ch. Pallinger}

% \begin{multicols}{2}


% \end{multicols}



\begin{WidePar}
\MarginBoxFilledRounded{Vorbereitung}{}
\bigskip
\begin{multicols}{2}
\begin{enumerate}
    \item Gewünschte Spatelgröße erfragen
    \item Gewünschte Tubusgröße erfragen
    \item Laryngoskop (Spatel und Griff) zusammenbauen
    \item Tubusverpackung aufschälen, Tubus aber in der Verpackung belassen   
    \item Dichtigkeit des Cuff prüfen: testweise aufblasen
    \item Führungsdraht (Mandrin) mit Silikonspray
    einsprühen\ACh{HART}{0.8.0}{Mandrin auch mit Silikon einsprühen} und
    einführen. Die Spitze muss bis zum Tubusende vorgeschoben vorgeschoben
    werden, darf aber nicht herausragen.
    \item Cuff mit Silikonspray einsprühen
    \item Stethoskop: Funktion prüfen
    \item Absaugung bereit machen:
    \begin{enumerate}
        \item Passenden Absaugkatheter wählen
        \item Verpackung teilweise aufschälen um Anschluss freizulegen, Katheter
        in der Verpackung belassen.
        \item Katheter anschließen.
        \item Funktionskontrolle: Gerät einschalten. Saugstärke?
        Batteriewarnung?
\item Absaugeinheit muss in Griffweite der Assistenten positioniert werden! 
    \end{enumerate}
    \item Beatmungsbeutel zusammenbauen und an \ce{O2}-Beriselungseinheit
    anschließen
    \item Restliches Material bereitlegen (z.\,B. in Nierentasse) 
\end{enumerate}
\end{multicols}
\end{WidePar}


% \Layout{2}{Assistenz}{
\begin{WidePar}
\MarginBoxFilledRounded{Assistenz}{}
\bigskip
\begin{multicols}{2}
\begin{enumerate}
  \item Überprüfen ob alle Materialien bereit sind.
  \item Sauerstoffberieselung einschalten und auf 15{\LPerMin*} einstellen,
  Reservoir des Beatmungsbeutels füllen.
  \item Durchführendem melden, dass alle Materialien bereit sind.
  \item \Z{Evtl. werden jetzt von einem Arzt Medikamente verabreicht, die
  die Intubation ermöglichen. Der Durchführende wird nun einige Zeit den
  Patienten mittels Beutel-/Masken-Beatmung beatmen (sofern nicht bereits
  geschehen).}
  \item Laryngoskop in die \EE{linke Hand} zureichen\footnote{Es ist egal, ob
  der Durchführende Rechts- oder Linkshänder ist, das Laryngoskop wird
  \E{immer in der linken Hand} gehalten.}.
  \item Auf Anweisung \emph{"`Kehlkopfdruck"', "`Krikoiddruck"'} auf den
  Kehlkopf drücken.
  \item Tubus in die rechte Hand zureichen.
  \item Auf Anweisung absaugen.
  \item Beatmungsmaske vom Beatmungsbeutel trennen.
  \item \Z{Durchführender intubiert und entfernt Führungsdraht (Mandrin)}
  \item Mit luftgefüllter Spritze Cuff aufblasen (\E{"`cuffen"'})
  \item Beatmungsbeutel an Tubus anschließen
  \item Notarzt Stethoskop in die Ohren klemmen.
  \item Es werden 3 Atemhübe verabreicht, wobei das
  Stethoskop über über dem Magen, linkem Lungenflügel und
  rechtem Lungenflügel positioniert wird. Es wird dabei überprüft ob der Tubus richtig liegt.
  \item Der Durchführende übernimmt die Beatmung.
  \item Einführen des Beißkeils oder des Guedel-Tubus
  \item Fixierung des Beißschutzes und des Tubus mittels Mullbinde
  \item \Z{Ggfs. Beatmungsgerät einstellen und
  anschließen}
\end{enumerate}
\end{multicols}
\end{WidePar}
% }{}{}{}{}

\TakeNotesMedium

\nocite{ZechaOpGehilfe1}
